% संपूर्ण मराठी ग्रंथातील प्रकरण 11: टॅब्लो
% हा संपादनयोग्य स्रोतखंड आहे. संकलनासाठी संपूर्ण स्रोतसंग्रहातील
% अक्षररूपे, पूर्वभाग, चिन्हव्याख्या आणि आकृती-स्रोत आवश्यक आहेत.
% मूळ: Open Logic Project; मजकूर आणि रूपांतर CC BY 4.0.
% यंत्रानुवाद, दुरुस्ती आणि तपासणी: OpenAI Codex —
% GPT-5.6 Sol आणि GPT-6 Sol, दोन्ही Ultra effort.
% दुरुस्ती-पुनरावलोकन: GPT-6.1 Sol, Ultra effort.
\chapter{टॅब्लो}\label{fol:tab::chap}\begingroup\setlength{\emergencystretch}{3em}\NewDocumentCommand{\sFmlaWide}{m m o}{\ensuremath{\IfNoValueTF{#3}{}{#3\,}\hbox to 1.2em{\ensuremath{#1}\hfil} #2}}
\begin{editorial}
हे प्रकरण चिन्हांकित सूत्रांवर आधारलेली विश्लेषणात्मक टॅब्लो पद्धत मांडते.
टॅब्लोशी सिद्धता-पद्धती म्हणून संबंधित सामग्री समाविष्ट किंवा
वगळण्यासाठी ``prfTab'' टॅग वापरा.
\end{editorial}
\section{नियम आणि टॅब्लो}\label{fol:tab:rul:sec}
टॅब्लो म्हणजे वाक्य संरचना मध्ये ज्या संभाव्य प्रकारे
सत्य अथवा असत्य असू शकते त्यांचा पद्धतशीर आढावा होय. टॅब्लोचे मूलभूत
घटक चिन्हांकित सूत्रे असतात: वाक्ये सोबत सत्यतामूल्याचे
``चिन्ह,'' म्हणजे $\True$ किंवा~$\False$. ही चिन्हांकित सूत्रे
(खाली वाढणाऱ्या) वृक्षात मांडली जातात.
\begin{defn}
  एक \emph{चिन्हांकित सूत्र} म्हणजे सत्यतामूल्य आणि वाक्य यांची जोडी,
  म्हणजे पुढीलपैकी एक:
  \[
  \sFmla{\True}{\varphi} \text{ किंवा } \sFmla{\False}{\varphi}.
  \]
\end{defn}

सहज अर्थानुसार, $\sFmla{\True}{\varphi}$ चे वाचन आपण ``$\varphi$ सत्य असू
शकते'' आणि $\sFmla{\False}{\varphi}$ चे वाचन ``$\varphi$ असत्य असू शकते''
(एखाद्या संरचना मध्ये) असे करू शकतो.

वृक्षातील प्रत्येक चिन्हांकित सूत्र एकतर \emph{गृहीतक} असते (ती
वृक्षाच्या अगदी वरच्या टोकाला नोंदवलेली असतात), किंवा तिच्या वर असलेल्या
चिन्हांकित सूत्राला अनेक निगमन नियमांपैकी एखादा नियम लावून ती मिळते.
पूर्ववर्ती सूत्राच्या प्रत्येक संभाव्य मुख्य संकारक साठी दोन नियम
असतात: चिन्ह~$\True$ असेल त्या स्थितीसाठी एक आणि चिन्ह~$\False$ असेल त्या
स्थितीसाठी एक. काही नियम वृक्षाची शाखा फोडू देतात, तर काही नियम शाखेत फक्त
चिन्हांकित सूत्रे जोडतात. एखादा नियम फक्त अगदी आधी येणाऱ्या
चिन्हांकित सूत्र वरच लागू करावा असे नाही; मुळापासून नियम लागू करण्याच्या
स्थानापर्यंतच्या शाखेतील कोणत्याही चिन्हांकित सूत्र वर तो लागू करता येतो
(आणि अनेकदा आधीच्या एखाद्यावर लावावाच लागतो).

एखाद्या शाखेत $\sFmla{\True}{\varphi}$ आणि $\sFmla{\False}{\varphi}$ ही दोन्ही
असतील तेव्हा ती शाखा \emph{बंद} असते. ज्याची प्रत्येक शाखा बंद असते तो
टॅब्लो बंद असतो. सहज अर्थानुसार, कोणतीही शाखा एक संयुक्त शक्यता
वर्णन करते; परंतु $\sFmla{\True}{\varphi}$ आणि $\sFmla{\False}{\varphi}$ एकाच वेळी
शक्य नाहीत. दुसऱ्या शब्दांत, एखादी शाखा बंद असेल तर तिने वर्णन केलेली शक्यता
बाद होते. विशेषतः याचा अर्थ असा की, बंद टॅब्लो हा
$\sFmla{\True}{\varphi}$ स्वरूपाचे प्रत्येक गृहीतक सत्य आणि
$\sFmla{\False}{\varphi}$ स्वरूपाचे प्रत्येक गृहीतक असत्य करणाऱ्या सर्व शक्यता
बाद करतो.

बंद टॅब्लो \emph{$\varphi$ साठीचा} म्हणजे ज्याचे
मूळ~$\sFmla{\False}{\varphi}$ आहे असा बंद टॅब्लो. असा बंद टॅब्लो
अस्तित्वात असल्यास, $\varphi$ असत्य असण्याच्या सर्व शक्यता बाद होतात; म्हणजेच
$\varphi$ प्रत्येक संरचना मध्ये सत्य असले पाहिजे.












\section{विधानीय नियम}\label{fol:tab:prl:sec}

\subsection{$\lnot$ साठीचे नियम}

\begin{defish}
\AxiomC{\sFmla{\True}{\lnot \varphi}}
\RightLabel{\TRule{\True}{\lnot}}
\UnaryInfC{\sFmla{\False}{\varphi}}
\DisplayProof
\hfill
\AxiomC{\sFmla{\False}{\lnot \varphi}}
\RightLabel{\TRule{\False}{\lnot}}
\UnaryInfC{\sFmla{\True}{\varphi}}
\DisplayProof
\end{defish}

\subsection{$\land$ साठीचे नियम}

\begin{defish}\noindent
\AxiomC{\sFmla{\True}{\varphi \land \psi}}
\RightLabel{\TRule{\True}{\land}}
\UnaryInfC{\sFmla{\True}{\varphi}}
\noLine
\UnaryInfC{\sFmla{\True}{\psi}}
\DisplayProof
\hfill
\AxiomC{\sFmla{\False}{\varphi \land \psi}}
\RightLabel{\TRule{\False}{\land}}
\UnaryInfC{$\sFmla{\False}{\varphi} \quad \mid \quad \sFmla{\False}{\psi}$}
\DisplayProof
\end{defish}

\subsection{$\lor$ साठीचे नियम}

\begin{defish}
\AxiomC{\sFmla{\True}{\varphi \lor \psi}}
\RightLabel{\TRule{\True}{\lor}}
\UnaryInfC{$\sFmla{\True}{\varphi} \quad \mid \quad \sFmla{\True}{\psi}$}
\DisplayProof
\hfill
\AxiomC{\sFmla{\False}{\varphi \lor \psi}}
\RightLabel{\TRule{\False}{\lor}}
\UnaryInfC{\sFmla{\False}{\varphi}}
\noLine
\UnaryInfC{\sFmla{\False}{\psi}}
\DisplayProof
\end{defish}

\subsection{$\lif$ साठीचे नियम}

\begin{defish}
\AxiomC{\sFmla{\True}{\varphi \lif \psi}}
\RightLabel{\TRule{\True}{\lif}}
\UnaryInfC{$\sFmla{\False}{\varphi} \quad \mid \quad \sFmla{\True}{\psi}$}
\DisplayProof
\hfill
\AxiomC{\sFmla{\False}{\varphi \lif \psi}}
\RightLabel{\TRule{\False}{\lif}}
\UnaryInfC{\sFmla{\True}{\varphi}}
\noLine
\UnaryInfC{\sFmla{\False}{\psi}}
\DisplayProof
\end{defish}

\subsection{कट नियम}

\begin{defish}
\AxiomC{}
\RightLabel{\Cut}
\UnaryInfC{$\sFmla{\True}{\varphi} \quad \mid \quad \sFmla{\False}{\varphi}$}
\DisplayProof
\end{defish}

\Cut{} नियम पूर्वीच्या चिन्हांकित सूत्र वर ``लावला'' जात नाही; त्याऐवजी
तो टॅब्लो मधील प्रत्येक शाखा दोन शाखांत फोडू देतो: एका शाखेत
$\sFmla{\True}{\varphi}$ आणि दुसरीत~$\sFmla{\False}{\varphi}$ असते. हा नियम आवश्यक
नाही---ज्या चिन्हांकित सूत्रांच्या संचासाठी बंद टॅब्लो आहे, त्या
संचासाठी \Cut{} न वापरणाराही एक बंद टॅब्लो असतो---परंतु त्यामुळे
टॅब्लो सोयीस्कर रीतीने एकत्र जोडता येतात.












\section{संख्यापकांचे नियम}\label{fol:tab:qrl:sec}

\subsection{$\lforall$ साठीचे नियम}

\begin{defish}
\AxiomC{\sFmla{\True}{\lforall[x][\varphi(x)]}}
\RightLabel{\TRule{\True}{\forall}}
\UnaryInfC{\sFmla{\True}{\varphi(t)}}
\DisplayProof
\hfill
\AxiomC{\sFmla{\False}{\lforall[x][\varphi(x)]}}
\RightLabel{\TRule{\False}{\lforall}}
\UnaryInfC{\sFmla{\False}{\varphi(a)}}
\DisplayProof
\end{defish}

\TRule{\True}{\lforall} मध्ये, $t$ हे बंद पद आहे (म्हणजे, चरे नसलेले
पद). \TRule{\False}{\lforall} मध्ये, $a$~हा असा स्थिरांक आहे की तो
\TRule{\False}{\lforall} नियमाच्या वरच्या शाखेत कोठेही येता कामा नये.
\TRule{\False}{\forall} अनुमानाचा $a$ हा \emph{आयगेन चर} आहे, असे आपण
म्हणतो.\footnote{वरील नियमात $a$ हा स्थिरांक असला तरी आपण
``आयगेन चर'' ही संज्ञा वापरतो. यामागे ऐतिहासिक कारणे आहेत.}

\subsection{$\lexists$ साठीचे नियम}

\begin{defish}
\AxiomC{\sFmla{\True}{\lexists[x][\varphi(x)]}}
\RightLabel{\TRule{\True}{\lexists}}
\UnaryInfC{\sFmla{\True}{\varphi(a)}}
\DisplayProof
\hfill
\AxiomC{\sFmla{\False}{\lexists[x][\varphi(x)]}}
\RightLabel{\TRule{\False}{\lexists}}
\UnaryInfC{\sFmla{\False}{\varphi(t)}}
\DisplayProof
\end{defish}

पुन्हा, $t$~हे बंद पद आहे आणि $a$~हा असा स्थिरांक आहे की तो
\TRule{\True}{\lexists} नियमाच्या वरच्या शाखेत येत नाही. आपण $a$ ला
\TRule{\True}{\lexists} अनुमानाचा \emph{आयगेन चर} म्हणतो.

\TRule{\False}{\lforall} किंवा \TRule{\True}{\lexists} अनुमानाच्या वरच्या
शाखेत आयगेन चर येत नसावा, या अटीला \emph{आयगेन चराची अट} म्हणतात.

\begin{explain}
ही प्रथा आठवा: $\varphi$ हे असे सूत्र असेल की त्यात
चल~$x$ मुक्त असेल, तर हे आपण~$\varphi(x)$ असे लिहून दर्शवतो.
त्याच संदर्भात, मग $\varphi(t)$ हा~$\Subst{\varphi}{t}{x}$ चा संक्षेप असतो.
त्यामुळे \TRule{\False}{\lexists} नियम पुढीलप्रमाणेही लिहिता येईल:
\begin{prooftree}
  \AxiomC{\sFmla{\False}{\lexists[x][\varphi]}}
  \RightLabel{\TRule{\False}{\lexists}}
  \UnaryInfC{\sFmla{\False}{\Subst{\varphi}{t}{x}}}
\end{prooftree}
लक्षात घ्या की~$\varphi$ मध्ये $t$ आधीच येऊ शकते; उदा., $\varphi$ हे
~$\Atom{\Obj P}{t,x}$ असू शकते. म्हणून,
$\sFmla{\False}{\lexists[x][\Atom{\Obj P}{t,x}]}$ पासून
$\sFmla{\False}{\Atom{\Obj P}{t,t}}$ निष्पन्न करणे हा
\TRule{\False}{\lexists} चा योग्य उपयोग आहे. परंतु
\TRule{\False}{\lforall} आणि~\TRule{\True}{\lexists} मधील आयगेन चराच्या
अटींनुसार स्थिरांक~$a$ हा~$\varphi$ मध्ये येता कामा नये. म्हणून,
$\sFmla{\False}{\lforall[x][\Atom{\Obj P}{a,x}]}$ पासून
$\sFmla{\False}{\Atom{\Obj P}{a,a}}$ हे
~$\TRule{\False}{\lforall}$ वापरून योग्य रीतीने निष्पन्न करता येत नाही.
\end{explain}

\begin{explain}
\TRule{\True}{\lforall} आणि \TRule{\False}{\lexists} मध्ये पद~$t$ बंद
असण्याव्यतिरिक्त त्यावर आणखी कोणतेही निर्बंध नाहीत. याउलट,
\TRule{\True}{\lexists} आणि \TRule{\False}{\lforall} नियमांमध्ये आयगेन
चराच्या अटीनुसार संबंधित अनुमानांच्या वरच्या शाखांमध्ये स्थिरांक~$a$
कोठेही येता कामा नये. ही पद्धत निर्दोष राहील याची खात्री करण्यासाठी ही अट
आवश्यक आहे. ही अट नसती, तर पुढील आकृती
$\lexists[x][\varphi(x)] \lif \lforall[x][\varphi(x)]$ साठीचा
बंद टॅब्लो ठरली असती:
\begin{center}
\begin{tableau}{}
  [\sFmla{\False}{\lexists[x][\varphi(x)] \lif \lforall[x][\varphi(x)]}, just=\TAss
    [\sFmla{\True}{\lexists[x][\varphi(x)]},
      just={\TRule{\False}{\lif}[1]}
      [\sFmla{\False}{\lforall[x][\varphi(x)]},
        just={\TRule{\False}{\lif}[1]}
        [\sFmla{\True}{\varphi(a)},
          just={\TRule{\True}{\lexists}[2]}
          [\sFmla{\False}{\varphi(a)},
            just={\TRule{\False}{\lforall}[3]}, close]
        ]
      ]
    ]
  ]
\end{tableau}
\end{center}
परंतु $\lexists[x][\varphi(x)] \lif
\lforall[x][\varphi(x)]$ हे वैध नाही.
\end{explain}












\section{टॅब्लो}\label{fol:tab:der:sec}

\begin{explain}
गृहीतक म्हणजे काय हे आपण सांगितले आहे आणि निगमन नियम दिले आहेत.
टॅब्लो यांपासून विगमनाधारित रीतीने निर्माण होतात: प्रत्येक
टॅब्लो ही एकतर एक किंवा अधिक गृहीतकांनी बनलेली एकच शाखा असते, किंवा
टॅब्लोच्या एखाद्या शाखेवर निगमन नियमांपैकी एक नियम लावल्याने ती मिळते.
\end{explain}

\begin{defn}[टॅब्लो]
$\sFmlaWide{S_1}{\varphi_1}$, \dots,\linebreak $\sFmlaWide{S_n}{\varphi_n}$ या गृहीतकांसाठीचा
टॅब्लो (येथे प्रत्येक $S_i$ हे $\True$ किंवा~$\False$ यांपैकी एक
आहे) म्हणजे पुढील अटी पूर्ण करणारा चिन्हांकित सूत्रांचा सांत वृक्ष होय:
\begin{enumerate}
\item वृक्षातील सर्वांत वरची $n$ चिन्हांकित सूत्रे म्हणजे एकाखाली एक
  असलेली\linebreak $\sFmlaWide{S_i}{\varphi_i}$ ही सूत्रे होत.
\item वृक्षातील गृहीतक नसलेली प्रत्येक चिन्हांकित सूत्र तिच्या वरच्या
  शाखेतील चिन्हांकित सूत्राला निगमन नियम योग्य रीतीने लावल्याने मिळते.
\end{enumerate}
टॅब्लोची एखादी शाखा $\sFmla{\True}{\varphi}$ आणि
~$\sFmla{\False}{\varphi}$ ही दोन्ही सूत्रे धारण करत असेल तर आणि तरच ती
\emph{बंद} असते; अन्यथा ती \emph{खुली} असते. ज्याची प्रत्येक शाखा बंद
असते तो टॅब्लो (त्याच्या गृहीतकांच्या संचासाठीचा)
\emph{बंद टॅब्लो} असतो. टॅब्लो बंद नसेल, म्हणजे त्यात किमान
एक खुली शाखा असेल, तर तो \emph{खुला} असतो.
\end{defn}

\begin{ex}
  गृहीतकांचा प्रत्येक संच स्वतःच टॅब्लो असतो, परंतु साधारणतः तो बंद
  नसतो. (स्पष्टपणे, तो बंद असेल तर गृहीतकांत $\sFmla{\True}{\varphi}$ आणि
  ~$\sFmla{\False}{\varphi}$ ही चिन्हांकित सूत्रांची जोडी आधीच असली पाहिजे.)

  एखाद्या टॅब्लो मधील चिन्हांकित सूत्र~$\varphi$ ला निगमन नियमांपैकी
  एक नियम लावून त्या टॅब्लोपासून (तो खुला असो वा बंद) नवा आणि अधिक मोठा
  टॅब्लो मिळवता येतो. नियम लागू केलेली~$\varphi$ ची आस्थिती ज्या ज्या शाखेत
  आहे, त्या प्रत्येक शाखेच्या शेवटी तो नियम एक किंवा अधिक
  चिन्हांकित सूत्रे जोडतो.

  उदाहरणार्थ, $\sFmla{\True}{\varphi \land \lnot \varphi}$ हे गृहीतक विचारात घ्या.
  केवळ त्या गृहीतकाने बनलेला (खुला) टॅब्लो पुढीलप्रमाणे आहे:
  \begin{center}
    \begin{tableau}{}
      [\sFmla{\True}{\varphi \land \lnot \varphi}, just=\TAss]
    \end{tableau}{}
  \end{center}
  त्या गृहीतकाला $\TRule{\True}{\land}$ नियम लावून त्यापासून नवा
  टॅब्लो मिळतो. हा नियम टॅब्लोमध्ये $\sFmla{\True}{\varphi}$ आणि
  $\sFmla{\True}{\lnot \varphi}$ या दोन नव्या ओळी जोडू देतो:
  \begin{center}
    \begin{tableau}{}
      [\sFmla{\True}{\varphi \land \lnot \varphi}, just=\TAss
        [\sFmla{\True}{\varphi}, just={\TRule{\True}{\land}[1]},
          [\sFmla{\True}{\lnot \varphi}, just={\TRule{\True}{\land}[1]}
          ]
        ]
      ]
    \end{tableau}{}
  \end{center}
  टॅब्लो लिहिताना आपण लावलेले नियम उजवीकडे नोंदवतो (उदा.,
  $\TRule{\True}{\land} 1$ याचा अर्थ त्या ओळीवरील चिन्हांकित सूत्र ही
  ओळ~$1$ वरील चिन्हांकित सूत्राला $\TRule{\True}{\land}$ नियम लावल्याने
  मिळाली आहे). या नव्या टॅब्लोमध्ये आता अधिक चिन्हांकित सूत्रे
  आहेत, परंतु त्यांपैकी फक्त एकीला ($\sFmla{\True}{\lnot \varphi}$) नियम लावता
  येतो (या स्थितीत $\TRule{\True}{\lnot}$ नियम). त्यामुळे पुढील बंद
  टॅब्लो मिळतो:
  \begin{center}
    \begin{tableau}{}
      [\sFmla{\True}{\varphi \land \lnot \varphi}, just=\TAss
        [\sFmla{\True}{\varphi}, just={\TRule{\True}{\land}[1]}
          [\sFmla{\True}{\lnot \varphi}, just={\TRule{\True}{\land}[1]}
            [\sFmla{\False}{\varphi}, just={\TRule{\True}{\lnot}[3]}, close]
          ]
        ]
      ]
    \end{tableau}{}
  \end{center}
\end{ex}












\section{टॅब्लोची उदाहरणे}\label{fol:tab:pro:sec}

\begin{ex}
$(\varphi \land \psi) \lif \varphi$ या वाक्य साठी बंद टॅब्लो शोधू या.

संबंधित गृहीतक टॅब्लोच्या सर्वांत वर लिहून आपण सुरुवात करतो.
\begin{oltableau}
  [\sFmla{\False}{(\varphi \land \psi) \lif \varphi},
    just = \TAss]
\end{oltableau}

एकच गृहीतक असल्यामुळे ज्याला नियम लावता येईल अशी एकच चिन्हांकित सूत्र
आहे. (प्रत्येक चिन्हांकित सूत्राला नेहमी जास्तीत जास्त एकच नियम लागू
करता येतो: तो म्हणजे वाक्याच्या संबंधित चिन्हाचा आणि
मुख्य संकारक चा नियम.) या स्थितीत याचा अर्थ असा की आपण
$\TRule{\False}{\lif}$ लावला पाहिजे.
\begin{oltableau}
  [\sFmla{\False}{(\varphi \land \psi) \lif \varphi},
    checked, just = \TAss
    [\sFmla{\True}{\varphi \land \psi},
      just={\TRule{\False}{\lif}[1]}
      [\sFmla{\False}{\varphi}, just={\TRule{\False}{\lif}[1]}]
    ]
  ]
\end{oltableau}
कोणत्या चिन्हांकित सूत्रे ना त्यांचे संबंधित नियम लावले आहेत याची नोंद
ठेवण्यासाठी आपण त्या वाक्याशेजारी खूण करतो. मात्र नियम सर्व खुल्या शाखांवर
लावला असेल \emph{तरच} खूण करा. एखाद्या चिन्हांकित सूत्राला प्रत्येक
खुल्या शाखेत संबंधित नियम लावल्यानंतर तिच्याकडे परत जाऊन तो नियम पुन्हा
लावावा लागणार नाही. येथे एकच शाखा असल्यामुळे नियम एकदाच लावावा लागतो.
(या खुणा केवळ टॅब्लो तयार करण्याची सोय आहेत आणि टॅब्लोच्या विन्यासाचा
अधिकृत भाग नाहीत, हे लक्षात घ्या.)

नियम लावता येईल अशी एक नवी चिन्हांकित सूत्र आहे: ओळ~$2$ वरील
$\sFmla{\True}{\varphi \land \psi}$. $\TRule{\True}{\land}$ नियम लावल्यावर
पुढील परिणाम मिळतो:
\begin{oltableau}
  [\sFmla{\False}{(\varphi \land \psi) \lif \varphi},
    checked, just = \TAss
    [\sFmla{\True}{\varphi \land \psi},
      just={\TRule{\False}{\lif}[1]}, checked
      [\sFmla{\False}{\varphi}, just={\TRule{\False}{\lif}[1]}
        [\sFmla{\True}{\varphi}, just={\TRule{\True}{\land}[2]}
          [\sFmla{\True}{\psi}, just={\TRule{\True}{\land}[2]}, close
          ]
        ]
      ]
    ]
  ]
\end{oltableau}
शाखेत आता $\sFmla{\True}{\varphi}$ (ओळ~$4$ वर) आणि
$\sFmla{\False}{\varphi}$ (ओळ~$3$ वर) ही दोन्ही असल्यामुळे शाखा बंद आहे.
ती एकमेव शाखा असल्यामुळे टॅब्लो बंद आहे. आपण
~$(\varphi \land \psi) \lif \varphi$ साठी बंद टॅब्लो शोधला आहे.
\end{ex}

\begin{ex}
आता $(\lnot \varphi \lor \psi) \lif (\varphi \lif \psi)$ साठी बंद टॅब्लो शोधू या.

संबंधित गृहीतकाने आपण सुरुवात करतो:
\begin{oltableau}
  [\sFmla{\False}{(\lnot \varphi \lor \psi) \lif
      (\varphi \lif \psi)}, just=\TAss]
\end{oltableau}
या टॅब्लो मधील एकमेव चिन्हांकित सूत्राचा मुख्य संकारक~$\lif$
आणि चिन्ह~$\False$ आहे; म्हणून तिला $\TRule{\False}{\lif}$ नियम लावल्यावर
पुढील परिणाम मिळतो:
\begin{oltableau}
  [\sFmla{\False}{(\lnot \varphi \lor \psi)
      \lif (\varphi \lif \psi)}, just=\TAss, checked
    [\sFmla{\True}{\lnot \varphi \lor \psi},
      just={\TRule{\False}{\lif}[1]}
      [\sFmla{\False}{(\varphi \lif \psi)},
        just={\TRule{\False}{\lif}[1]}
      ]
    ]
  ]
\end{oltableau}
आता ओळ~$2$ ला~$\TRule{\True}{\lor}$ लावायचा की ओळ~$3$ ला
$\TRule{\False}{\lif}$ लावायचा, असा पर्याय आपल्यापुढे आहे. प्रत्येक
चिन्हांकित सूत्राला तिचा संबंधित नियम प्रत्येक शाखेत लावला, तर आपण कोणता
क्रम निवडतो याने प्रत्यक्षात फरक पडत नाही. म्हणून पहिला पर्याय निवडू या.
$\TRule{\True}{\lor}$ नियमामुळे टॅब्लोची शाखा फुटू शकते आणि नियमाचे
दोन निष्कर्ष त्या दोन नव्या शाखांना जोडलेल्या नव्या चिन्हांकित सूत्रे
असतील. त्यामुळे पुढील परिणाम मिळतो:
\begin{oltableau}
  [\sFmla{\False}{(\lnot \varphi \lor \psi) \lif
      (\varphi \lif \psi)}, just=\TAss, checked
    [\sFmla{\True}{\lnot \varphi \lor \psi},
      just={\TRule{\False}{\lif}[1]}, checked
      [\sFmla{\False}{(\varphi \lif \psi)},
        just={\TRule{\False}{\lif}[1]}
        [\sFmla{\True}{\lnot \varphi}, just={\TRule{\True}{\lor}[2]}]
        [\sFmla{\True}{\psi}, just={\TRule{\True}{\lor}[2]}]
      ]
    ]
  ]
\end{oltableau}
ओळ~$3$ ला आपण अद्याप $\TRule{\False}{\lif}$ नियम लावलेला नाही; तो आता
लावू या. वेळ वाचवण्यासाठी तो दोन्ही शाखांवर लावू. प्रत्येक खुल्या शाखेत
संबंधित नियम लावला असेल तरच चिन्हांकित सूत्र शेजारी खूण करतो, हे आठवा.
म्हणून नियम ज्याला लागू होतो ती चिन्हांकित सूत्र असलेल्या प्रत्येक
शाखेच्या शेवटी तो नियम लावणे उचित ठरते. त्यामुळे विविध शाखांत खाली जाऊन त्या
चिन्हांकित सूत्र कडे पुन्हा परतावे लागत नाही.
\begin{oltableau}
  [\sFmla{\False}{(\lnot \varphi \lor \psi) \lif
      (\varphi \lif \psi)}, just=\TAss, checked
    [\sFmla{\True}{\lnot \varphi \lor \psi},
      just={\TRule{\False}{\lif}[1]}, checked
      [\sFmla{\False}{(\varphi \lif \psi)},
        just={\TRule{\False}{\lif}[1]}, checked
        [\sFmla{\True}{\lnot \varphi}, just={\TRule{\True}{\lor}[2]}
          [\sFmla{\True}{\varphi}, just={\TRule{\False}{\lif}[3]}
            [\sFmla{\False}{\psi}, just={\TRule{\False}{\lif}[3]}]
          ]
        ]
        [\sFmla{\True}{\psi}, just={\TRule{\True}{\lor}[2]}
          [\sFmla{\True}{\varphi}, just={\TRule{\False}{\lif}[3]}
            [\sFmla{\False}{\psi}, just={\TRule{\False}{\lif}[3]}, close]
          ]
        ]
      ]
    ]
  ]
\end{oltableau}
उजवी शाखा आता बंद आहे. डाव्या शाखेत आपण अजूनही ओळ~$4$ ला
$\TRule{\True}{\lnot}$ नियम लावू शकतो. त्यामुळे
~$\sFmla{\False}{\varphi}$ मिळते आणि डावी शाखा बंद होते:
\begin{oltableau}
  [\sFmla{\False}{(\lnot \varphi \lor \psi) \lif
      (\varphi \lif \psi)}, just=\TAss, checked
    [\sFmla{\True}{\lnot \varphi \lor \psi},
      just={\TRule{\False}{\lif}[1]}, checked
      [\sFmla{\False}{(\varphi \lif \psi)},
        just={\TRule{\False}{\lif}[1]}, checked
        [\sFmla{\True}{\lnot \varphi}, just={\TRule{\True}{\lor}[2]}
          [\sFmla{\True}{\varphi}, just={\TRule{\False}{\lif}[3]}
            [\sFmla{\False}{\psi}, just={\TRule{\False}{\lif}[3]}
              [\sFmla{\False}{\varphi},
                just={\TRule{\True}{\lnot}[4]}, close]
            ]
          ]
        ]
        [\sFmla{\True}{\psi}, just={\TRule{\True}{\lor}[2]}
          [\sFmla{\True}{\varphi}, just={\TRule{\False}{\lif}[3]}
            [\sFmla{\False}{\psi},
              just={\TRule{\False}{\lif}[3]}, close]
          ]
        ]
      ]
    ]
  ]
\end{oltableau}
\end{ex}

\begin{ex}
कितीही चिन्हांकित सूत्रे गृहीतके म्हणून घेऊन त्यांच्यासाठी टॅब्लो
देता येतात. शाखा फोडू देणारे एकाहून अधिक नियम लावणेही अनेकदा आवश्यक असते;
आणि सर्वसाधारणपणे टॅब्लोमध्ये कितीही शाखा असू शकतात. उदाहरणार्थ,
पुढील सूत्रांसाठीचा टॅब्लो विचारात घ्या:
$\{\sFmla{\True}{\varphi \lor (\psi \land \chi)}, \sFmla{\False}{(\varphi \lor \psi) \land (\varphi
  \lor \chi)}\}$. पहिल्या गृहीतकाला $\TRule{\True}{\lor}$ लावून आपण
सुरुवात करतो:
\begin{oltableau}
  [\sFmla{\True}{\varphi \lor (\psi \land \formula{C})},
    just=\TAss, checked
    [\sFmla{\False}{(\varphi \lor \psi) \land
        (\varphi \lor \formula{C})}, just=\TAss
      [\sFmla{\True}{\varphi}, just={\TRule{\True}{\lor}[1]}]
      [\sFmla{\True}{\psi \land \formula{C}},
        just={\TRule{\True}{\lor}[1]}]
    ]
  ]
\end{oltableau}
आता ओळ~$2$ ला $\TRule{\False}{\land}$ नियम लावता येतो. दोन्ही शाखांवर
तो एकाच वेळी लावल्यामुळे आपण ओळ~$2$ वर खूण करू शकतो:
\begin{oltableau}
  [\sFmla{\True}{\varphi \lor (\psi \land \formula{C})},
    just=\TAss, checked
    [\sFmla{\False}{(\varphi \lor \psi) \land
        (\varphi \lor \formula{C})}, just=\TAss, checked
      [\sFmla{\True}{\varphi}, just={\TRule{\True}{\lor}[1]}
        [\sFmla{\False}{\varphi \lor \psi},
          just={\TRule{\False}{\land}[2]}]
        [\sFmla{\False}{\varphi \lor \formula{C}},
          just={\TRule{\False}{\land}[2]}]
      ]
      [\sFmla{\True}{\psi \land \formula{C}},
        just={\TRule{\True}{\lor}[1]}
        [\sFmla{\False}{\varphi \lor \psi},
          just={\TRule{\False}{\land}[2]}]
        [\sFmla{\False}{\varphi \lor \formula{C}},
          just={\TRule{\False}{\land}[2]}]
      ]
    ]
  ]
\end{oltableau}
आता $\varphi \lor \psi$ असलेल्या सर्व शाखांवर $\TRule{\False}{\lor}$ लावता
येतो:
\begin{oltableau}
  [\sFmla{\True}{\varphi \lor (\psi \land \formula{C})},
    just=\TAss, checked
    [\sFmla{\False}{(\varphi \lor \psi) \land
        (\varphi \lor \formula{C})}, just=\TAss, checked
      [\sFmla{\True}{\varphi}, just={\TRule{\True}{\lor}[1]}
        [\sFmla{\False}{\varphi \lor \psi},
          just={\TRule{\False}{\land}[2]}, checked
          [\sFmla{\False}{\varphi}, just={\TRule{\False}{\lor}[4]}
            [\sFmla{\False}{\psi},
              just={\TRule{\False}{\lor}[4]}, close]
          ]
        ]
        [\sFmla{\False}{\varphi \lor \formula{C}},
          just={\TRule{\False}{\land}[2]}]
      ]
      [\sFmla{\True}{\psi \land \formula{C}},
        just={\TRule{\True}{\lor}[1]}
        [\sFmla{\False}{\varphi \lor \psi},
          just={\TRule{\False}{\land}[2]}, checked
          [\sFmla{\False}{\varphi}, just={\TRule{\False}{\lor}[4]}
            [\sFmla{\False}{\psi},
              just={\TRule{\False}{\lor}[4]}]
          ]
        ]
        [\sFmla{\False}{\varphi \lor \formula{C}},
          just={\TRule{\False}{\land}[2]}]
      ]
    ]
  ]
\end{oltableau}
सर्वांत डावी शाखा आता बंद आहे. आता $\varphi \lor \chi$ ला
$\TRule{\False}{\lor}$ लावू या:
\begin{oltableau}
  [\sFmla{\True}{\varphi \lor (\psi \land \formula{C})},
    just=\TAss, checked
    [\sFmla{\False}{(\varphi \lor \psi) \land
        (\varphi \lor \formula{C})}, just=\TAss, checked
      [\sFmla{\True}{\varphi}, just={\TRule{\True}{\lor}[1]}
        [\sFmla{\False}{\varphi \lor \psi},
          just={\TRule{\False}{\land}[2]}, checked
          [\sFmla{\False}{\varphi},
            just={\TRule{\False}{\lor}[4]}
            [\sFmla{\False}{\psi},
              just={\TRule{\False}{\lor}[4]}, close]
          ]
        ]
        [\sFmla{\False}{\varphi \lor \formula{C}},
          just={\TRule{\False}{\land}[2]}, checked
          [\sFmla{\False}{\varphi},
            just={\TRule{\False}{\lor}[4]}, move by=2
            [\sFmla{\False}{\formula{C}},
              just={\TRule{\False}{\lor}[4]}, close]
          ]
        ]
      ]
      [\sFmla{\True}{\psi \land \formula{C}},
        just={\TRule{\True}{\lor}[1]}
        [\sFmla{\False}{\varphi \lor \psi},
          just={\TRule{\False}{\land}[2]}, checked
          [\sFmla{\False}{\varphi},
            just={\TRule{\False}{\lor}[4]}
            [\sFmla{\False}{\psi},
              just={\TRule{\False}{\lor}[4]}
            ]
          ]
        ]
        [\sFmla{\False}{\varphi \lor \formula{C}},
          just={\TRule{\False}{\land}[2]}, checked
          [\sFmla{\False}{\varphi},
            just={\TRule{\False}{\lor}[4]}, move by=2
            [\sFmla{\False}{\formula{C}},
              just={\TRule{\False}{\lor}[4]}]
          ]
        ]
      ]
    ]
  ]
\end{oltableau}
स्पष्टतेसाठी $\TRule{\False}{\lor}$ दुसऱ्यांदा लावल्याचा परिणाम आपण खाली
हलवला आहे, हे लक्षात घ्या. येथे समर्थन समान असते, त्यामुळे तसे करण्याची
गरज नव्हती.

दोन शाखा खुल्या राहतात आणि ओळ~$3$ वरील $\sFmla{\True}{\psi \land \chi}$
वर खूण झालेली नाही. तिला $\TRule{\True}{\land}$ लावून पुढील बंद
टॅब्लो मिळतो:
\begin{oltableau}
  [\sFmla{\True}{\varphi \lor (\psi \land \formula{C})},
    just=\TAss, checked
    [\sFmla{\False}{(\varphi \lor \psi) \land
        (\varphi \lor \formula{C})}, just=\TAss, checked
      [\sFmla{\True}{\varphi}, just={\TRule{\True}{\lor}[1]}
        [\sFmla{\False}{\varphi \lor \psi},
          just={\TRule{\False}{\land}[2]}, checked
          [\sFmla{\False}{\varphi},
            just={\TRule{\False}{\lor}[4]}
            [\sFmla{\False}{\psi},
              just={\TRule{\False}{\lor}[4]}, close]
          ]
        ]
        [\sFmla{\False}{\varphi \lor \formula{C}},
          just={\TRule{\False}{\land}[2]}, checked
          [\sFmla{\False}{\varphi},
            just={\TRule{\False}{\lor}[4]}
            [\sFmla{\False}{\formula{C}},
              just={\TRule{\False}{\lor}[4]}, close]
          ]
        ]
      ]
      [\sFmla{\True}{\psi \land \formula{C}},
        just={\TRule{\True}{\lor}[1]}, checked
        [\sFmla{\False}{\varphi \lor \psi},
          just={\TRule{\False}{\land}[2]}, checked
          [\sFmla{\False}{\varphi},
            just={\TRule{\False}{\lor}[4]}
            [\sFmla{\False}{\psi},
              just={\TRule{\False}{\lor}[4]}
              [\sFmla{\True}{\psi},
                just={\TRule{\True}{\land}[3]}
                [\sFmla{\True}{\formula{C}},
                  just={\TRule{\True}{\land}[3]},close
                ]
              ]
            ]
          ]
        ]
        [\sFmla{\False}{\varphi \lor \formula{C}},
          just={\TRule{\False}{\land}[2]}, checked
          [\sFmla{\False}{\varphi},
            just={\TRule{\False}{\lor}[4]}
            [\sFmla{\False}{\formula{C}},
              just={\TRule{\False}{\lor}[4]}
              [\sFmla{\True}{\psi},
                just={\TRule{\True}{\land}[3]}
                [\sFmla{\True}{\formula{C}},
                  just={\TRule{\True}{\land}[3]},close
                ]
              ]
            ]
          ]
        ]
      ]
    ]
  ]
\end{oltableau}

तुलनेसाठी, त्याच गृहीतकांच्या संचासाठी नियम वेगळ्या क्रमाने लावलेला बंद
टॅब्लो पुढे दिला आहे:
\begin{oltableau}
  [\sFmla{\True}{\varphi \lor (\psi \land \formula{C})},
    just=\TAss, checked
    [\sFmla{\False}{(\varphi \lor \psi) \land
        (\varphi \lor \formula{C})}, just=\TAss, checked
      [\sFmla{\False}{\varphi \lor \psi},
        just={\TRule{\False}{\land}[2]}, checked
        [\sFmla{\False}{\varphi},
          just={\TRule{\False}{\lor}[3]}
          [\sFmla{\False}{\psi},
            just={\TRule{\False}{\lor}[3]}
            [\sFmla{\True}{\varphi},
              just={\TRule{\True}{\lor}[1]},close
            ]
            [\sFmla{\True}{\psi \land \formula{C}},
              just={\TRule{\True}{\lor}[1]}, checked
              [\sFmla{\True}{\psi},
                just={\TRule{\True}{\land}[6]}
                [\sFmla{\True}{\formula{C}},
                  just={\TRule{\True}{\land}[6]},close
                ]
              ]
            ]
          ]
        ]
      ]
      [\sFmla{\False}{\varphi \lor \formula{C}},
        just={\TRule{\False}{\land}[2]}, checked
        [\sFmla{\False}{\varphi},
          just={\TRule{\False}{\lor}[3]}
          [\sFmla{\False}{\formula{C}},
            just={\TRule{\False}{\lor}[3]}
            [\sFmla{\True}{\varphi},
              just={\TRule{\True}{\lor}[1]},close
            ]
            [\sFmla{\True}{\psi \land \formula{C}},
              just={\TRule{\True}{\lor}[1]}, checked
              [\sFmla{\True}{\psi},
                just={\TRule{\True}{\land}[6]}
                [\sFmla{\True}{\formula{C}},
                  just={\TRule{\True}{\land}[6]},close
                ]
              ]
            ]
          ]
        ]
      ]
    ]
  ]
\end{oltableau}
\end{ex}

\begin{prob}
पुढील सूत्रांचे बंद टॅब्लो द्या:
\begin{enumerate}
\item $\sFmla{\True}{\varphi \land (\psi \land \chi)}, \sFmla{\False}{(\varphi \land \psi) \land \chi}$.
\item $\sFmla{\True}{\varphi \lor (\psi \lor \chi)}, \sFmla{\False}{(\varphi \lor \psi) \lor \chi}$.
\item $\sFmla{\True}{\varphi \lif (\psi \lif \chi)}, \sFmla{\False}{\psi \lif (\varphi \lif \chi)}$.
\item $\sFmla{\True}{\varphi}, \sFmla{\False}{\lnot\lnot \varphi}$.
\end{enumerate}
\end{prob}

\begin{prob}
पुढील सूत्रांचे बंद टॅब्लो द्या:
\begin{enumerate}
\item $\sFmla{\True}{(\varphi \lor \psi) \lif \chi}, \sFmla{\False}{\varphi \lif \chi}$.
\item $\sFmla{\True}{(\varphi \lif \chi) \land (\psi \lif \chi)}, \sFmla{\False}{(\varphi \lor \psi) \lif \chi}$.
\item $\sFmla{\False}{\lnot(\varphi \land \lnot \varphi)}$.
\item $\sFmla{\True}{\psi \lif \varphi}, \sFmla{\False}{\lnot \varphi \lif \lnot \psi}$.
\item $\sFmla{\False}{(\varphi \lif \lnot \varphi) \lif \lnot \varphi}$.
\item $\sFmla{\False}{\lnot(\varphi \lif \psi) \lif \lnot \psi}$.
\item $\sFmla{\True}{\varphi \lif \chi}, \sFmla{\False}{\lnot (\varphi \land \lnot \chi)}$.
\item $\sFmla{\True}{\varphi \land \lnot \chi}, \sFmla{\False}{\lnot (\varphi \lif \chi)}$.
\item $\sFmla{\True}{\varphi \lor \psi, \lnot \psi}, \sFmla{\False}{\varphi}$.
\item $\sFmla{\True}{\lnot \varphi \lor \lnot \psi}, \sFmla{\False}{\lnot(\varphi \land \psi)}$.
\item $\sFmla{\False}{(\lnot \varphi \land \lnot \psi) \lif\lnot(\varphi \lor \psi)}$.
\item $\sFmla{\False}{\lnot(\varphi \lor \psi) \lif (\lnot \varphi \land \lnot \psi)}$.
\end{enumerate}
\end{prob}

\begin{prob}
पुढील सूत्रांचे बंद टॅब्लो द्या:
\begin{enumerate}
\item $\sFmla{\True}{\lnot(\varphi \lif \psi)}, \sFmla{\False}{\varphi}$.
\item $\sFmla{\True}{\lnot(\varphi \land \psi)}, \sFmla{\False}{\lnot \varphi \lor \lnot \psi}$.
\item $\sFmla{\True}{\varphi \lif \psi}, \sFmla{\False}{\lnot \varphi \lor \psi}$.
\item $\sFmla{\False}{\lnot \lnot \varphi \lif \varphi}$.
\item $\sFmla{\True}{\varphi \lif \psi}, \sFmla{\True}{\lnot \varphi \lif \psi}, \sFmla{\False}{\psi}$.
\item $\sFmla{\True}{(\varphi \land \psi) \lif \chi}, \sFmla{\False}{(\varphi \lif \chi) \lor (\psi \lif \chi)}$.
\item $\sFmla{\True}{(\varphi \lif \psi) \lif \varphi}, \sFmla{\False}{\varphi}$.
\item $\sFmla{\False}{(\varphi \lif \psi) \lor (\psi \lif \chi)}$.
\end{enumerate}
\end{prob}












\section{संख्यापकांसह टॅब्लो}\label{fol:tab:prq:sec}

\begin{ex}
संख्यापक हाताळताना आयगेन चराच्या अटीचा भंग होणार नाही याची खात्री करावी
लागते; त्यासाठी कधीकधी काही अनुमाने करण्याचा क्रम बदलून पाहावा लागतो.
सामान्यतः आयगेन चराच्या अटीच्या अधीन असलेले नियम आधी हाताळण्याचा प्रयत्न
उपयोगी ठरतो (पूर्ण टॅब्लोमध्ये ते अधिक वर असतील).

$\lexists[x][\lnot \varphi(x)] \lif \lnot \lforall[x][\varphi(x)]$ या
वाक्य साठी टॅब्लो कसा द्यायचा ते पाहू. नेहमीप्रमाणे संबंधित
गृहीतक नोंदवून आपण सुरुवात करतो:
\begin{oltableau}
[\sFmla{\False}{\lexists[x][\lnot \varphi(x)]\lif \lnot
    \lforall[x][\varphi(x)]}, just=\TAss]
\end{oltableau}
मुख्य संकारक $\lif$ असल्यामुळे आपण $\TRule{\False}{\lif}$ लावतो:
\begin{oltableau}
[\sFmla{\False}{\lexists[x][\lnot \varphi(x)]\lif \lnot
    \lforall[x][\varphi(x)]}, just=\TAss, checked
  [\sFmla{\True}{\lexists[x][\lnot \varphi(x)]},
    just={\TRule{\False}{\lif}[1]}
    [\sFmla{\False}{\lnot\lforall[x][\varphi(x)]},
      just={\TRule{\False}{\lif}[1]}]
  ]
]
\end{oltableau}
यानंतर ओळ~$2$ हाताळायची आहे. आपण $\TRule{\True}{\lexists}$ वापरतो.
त्यासाठी नवा स्थिरांक आवश्यक आहे; अद्याप कोणतेही स्थिरांक आलेले
नसल्यामुळे कोणताही एक, समजा~$a$, निवडता येतो.
\begin{oltableau}
[\sFmla{\False}{\lexists[x][\lnot \varphi(x)]\lif \lnot
    \lforall[x][\varphi(x)]}, just=\TAss, checked
  [\sFmla{\True}{\lexists[x][\lnot \varphi(x)]},
    just={\TRule{\False}{\lif}[1]}, checked
    [\sFmla{\False}{\lnot\lforall[x][\varphi(x)]},
      just={\TRule{\False}{\lif}[1]}
      [\sFmla{\True}{\lnot\varphi(a)}, just={\TRule{\True}{\lexists}[2]}]
    ]
  ]
]
\end{oltableau}
आता ओळ~$3$ ला $\TRule{\False}{\lnot}$ लावतो:
\begin{oltableau}
[\sFmla{\False}{\lexists[x][\lnot \varphi(x)]\lif \lnot
    \lforall[x][\varphi(x)]}, just=\TAss, checked
  [\sFmla{\True}{\lexists[x][\lnot \varphi(x)]},
    just={\TRule{\False}{\lif}[1]}, checked
    [\sFmla{\False}{\lnot\lforall[x][\varphi(x)]},
      just={\TRule{\False}{\lif}[1]}, checked
      [\sFmla{\True}{\lnot\varphi(a)}, just={\TRule{\True}{\lexists}[2]}
        [\sFmla{\True}{\lforall[x][\varphi(x)]},
          just = {\TRule{\False}{\lnot}[3]}
        ]
      ]
    ]
  ]
]
\end{oltableau}
ओळ~$4$ ला $\TRule{\True}{\lnot}$ आणि त्यानंतर ओळ~$5$ ला
$\TRule{\True}{\lforall}$ लावल्यावर बंद टॅब्लो मिळतो.
\begin{oltableau}
[\sFmla{\False}{\lexists[x][\lnot \varphi(x)]\lif \lnot
    \lforall[x][\varphi(x)]}, just=\TAss, checked
  [\sFmla{\True}{\lexists[x][\lnot \varphi(x)]},
    just={\TRule{\False}{\lif}[1]}, checked
    [\sFmla{\False}{\lnot\lforall[x][\varphi(x)]},
      just={\TRule{\False}{\lif}[1]}, checked
      [\sFmla{\True}{\lnot\varphi(a)}, just={\TRule{\True}{\lexists}[2]}
        [\sFmla{\True}{\lforall[x][\varphi(x)]},
          just = {\TRule{\False}{\lnot}[3]}
          [\sFmla{\False}{\varphi(a)}, just={\TRule{\True}{\lnot}[4]}
            [\sFmla{\True}{\varphi(a)},
              just={\TRule{\True}{\lforall}[5]}, close
            ]
          ]
        ]
      ]
    ]
  ]
]
\end{oltableau}
\end{ex}

\begin{ex}
पुढील संचासाठी टॅब्लो कसा द्यायचा ते पाहू:
\[\sFmla{\False}{\lexists[x][\chi(x,b)]},
\sFmla{\True}{\lexists[x][(\varphi(x) \land \psi(x))]},
\sFmla{\True}{\lforall[x][(\psi(x) \lif \chi(x,b))]}.
\]
नेहमीप्रमाणे गृहीतकांनी सुरुवात करतो:
\begin{oltableau}
  [\sFmla{\False}{\lexists[x][\formula{C}(x,b)]}, just=\TAss
    [\sFmla{\True}{\lexists[x][(\varphi(x) \land \psi(x))]},
      just=\TAss
      [\sFmla{\True}{\lforall[x][(\psi(x) \lif \formula{C}(x,b))]},
        just=\TAss]
    ]
  ]
\end{oltableau}
आयगेन चराची अट असलेला नियम आपण नेहमी आधी लावला पाहिजे; येथे तो म्हणजे
ओळ~$2$ ला $\TRule{\True}{\lexists}$. गृहीतकांत स्थिरांक~$b$ असल्यामुळे
वेगळा अचर वापरावा लागतो; पुन्हा~$a$ निवडू या.
\begin{oltableau}
  [\sFmla{\False}{\lexists[x][\formula{C}(x,b)]}, just=\TAss
    [\sFmla{\True}{\lexists[x][(\varphi(x) \land \psi(x))]},
      just=\TAss, checked
      [\sFmla{\True}{\lforall[x][(\psi(x) \lif \formula{C}(x,b))]},
        just=\TAss
        [\sFmla{\True}{\varphi(a) \land \psi(a)},
          just={\TRule{\True}{\lexists}[2]}]
      ]
    ]
  ]
\end{oltableau}
आता ओळ~$1$ ला $\TRule{\False}{\lexists}$ किंवा ओळ~$3$ ला
$\TRule{\True}{\lforall}$ लावल्यास, $x$ च्या जागी कोणते बंद पद~$t$ ठेवायचे
ते ठरवावे लागते. या नियमांना आयगेन चराची अट नसल्यामुळे हवे ते कोणतेही बंद
पद निवडता येते. काही स्थितींमध्ये वेगवेगळी~$t$ वापरून नियम अनेकदा लावावा
लागू शकतो. मात्र सर्वसाधारणपणे टॅब्लोमध्ये आधीच आलेल्या पदांपैकी
एक---येथे $a$ किंवा~$b$---निवडणे फायदेशीर ठरते; आणि येथे $a$ मुळे बंद
शाखा मिळण्याची शक्यता अधिक आहे असा अंदाज करता येतो.
\begin{oltableau}
  [\sFmla{\False}{\lexists[x][\formula{C}(x,b)]}, just=\TAss
    [\sFmla{\True}{\lexists[x][(\varphi(x) \land \psi(x))]},
      just=\TAss, checked
      [\sFmla{\True}{\lforall[x][(\psi(x) \lif \formula{C}(x,b))]}, just=\TAss
        [\sFmla{\True}{\varphi(a) \land \psi(a)}, just={\TRule{\True}{\lexists}[2]}
          [\sFmla{\False}{\formula{C}(a,b)}, just={\TRule{\False}{\lexists}[1]}
            [\sFmla{\True}{\psi(a) \lif \formula{C}(a,b)},
              just={\TRule{\True}{\lforall}[3]}
            ]
          ]
        ]
      ]
    ]
  ]
\end{oltableau}
ओळी $1$ आणि~$3$ वरील चिन्हांकित सूत्रे पुन्हा वापराव्या लागू शकतात,
म्हणून त्यांवर खूण करत नाही. आता ओळ~$4$ ला $\TRule{\True}{\land}$ लावा:
\begin{oltableau}
  [\sFmla{\False}{\lexists[x][\formula{C}(x,b)]}, just=\TAss
    [\sFmla{\True}{\lexists[x][(\varphi(x) \land \psi(x))]},
      just=\TAss, checked
      [\sFmla{\True}{\lforall[x][(\psi(x) \lif \formula{C}(x,b))]}, just=\TAss
        [\sFmla{\True}{\varphi(a) \land \psi(a)},
          just={\TRule{\True}{\lexists}[2]}, checked
          [\sFmla{\False}{\formula{C}(a,b)}, just={\TRule{\False}{\lexists}[1]}
            [\sFmla{\True}{\psi(a) \lif \formula{C}(a,b)},
              just={\TRule{\True}{\lforall}[3]}
              [\sFmla{\True}{\varphi(a)}, just={\TRule{\True}{\land}[4]}
                [\sFmla{\True}{\psi(a)}, just={\TRule{\True}{\land}[4]}
                ]
              ]
            ]
          ]
        ]
      ]
    ]
  ]
\end{oltableau}
आता ओळ~$6$ ला $\TRule{\True}{\lif}$ लावल्यास टॅब्लो बंद होतो:
\begin{oltableau}
  [\sFmla{\False}{\lexists[x][\formula{C}(x,b)]}, just=\TAss
    [\sFmla{\True}{\lexists[x][(\varphi(x) \land \psi(x))]},
      just=\TAss, checked
      [\sFmla{\True}{\lforall[x][(\psi(x) \lif \formula{C}(x,b))]},
        just=\TAss
        [\sFmla{\True}{\varphi(a) \land \psi(a)},
          just={\TRule{\True}{\lexists}[2]}, checked
          [\sFmla{\False}{\formula{C}(a,b)},
            just={\TRule{\False}{\lexists}[1]}
            [\sFmla{\True}{\psi(a) \lif \formula{C}(a,b)},
              just={\TRule{\True}{\lforall}[3]}, checked
              [\sFmla{\True}{\varphi(a)},
                just={\TRule{\True}{\land}[4]}
                [\sFmla{\True}{\psi(a)},
                  just={\TRule{\True}{\land}[4]}
                  [\sFmla{\False}{\psi(a)},
                    just={\TRule{\True}{\lif}[6]},close
                  ]
                  [\sFmla{\True}{\formula{C}(a,b)},
                    just={\TRule{\True}{\lif}[6]},close
                  ]
                ]
              ]
            ]
          ]
        ]
      ]
    ]
  ]
\end{oltableau}


\end{ex}

\begin{ex}
पुढील संचासाठी आपण टॅब्लो तयार करतो:
\[\sFmla{\True}{\lforall[x][\varphi(x)]}, \sFmla{\True}{\lforall[x][\varphi(x)]
  \lif \lexists[y][\psi(y)]}, \sFmla{\True}{\lnot\lexists[y][\psi(y)]}.
\]
नेहमीप्रमाणे गृहीतके लिहून सुरुवात करतो:
\begin{oltableau}
  [\sFmla{\True}{\lforall[x][\varphi(x)]}, just=\TAss
    [\sFmla{\True}{\lforall[x][\varphi(x)] \lif
        \lexists[y][\psi(y)]}, just=\TAss
      [\sFmla{\True}{\lnot\lexists[y][\psi(y)]}, just=\TAss
      ]
    ]
  ]
\end{oltableau}
ओळ~$3$ ला $\TRule{\True}{\lnot}$ नियम लावून सुरुवात करतो. ``आयगेन
चराची अट असलेले नियम नेहमी आधी लावा'' या नियमाचा एक परिणाम असा की
``आयगेन चराची अट नसलेले संख्यापक-नियम आवश्यक होईपर्यंत पुढे ढकला.''
शाखा फोडणारे नियमही पुढे ढकला.
\begin{oltableau}
  [\sFmla{\True}{\lforall[x][\varphi(x)]}, just=\TAss
    [\sFmla{\True}{\lforall[x][\varphi(x)] \lif
        \lexists[y][\psi(y)]}, just=\TAss
      [\sFmla{\True}{\lnot\lexists[y][\psi(y)]}, just=\TAss, checked
        [\sFmla{\False}{\lexists[y][\psi(y)]}, just={\TRule{\True}{\lnot}[3]}]
      ]
    ]
  ]
\end{oltableau}
नव्या ओळ~$4$ साठी $\TRule{\False}{\lexists}$ आवश्यक आहे आणि हा आयगेन
चराची अट नसलेला संख्यापक-नियम आहे. म्हणून तो पुढे ढकलून त्याऐवजी ओळ~$2$
वर $\TRule{\True}{\lif}$ वापरतो.
\begin{oltableau}
  [\sFmla{\True}{\lforall[x][\varphi(x)]}, just=\TAss
    [\sFmla{\True}{\lforall[x][\varphi(x)] \lif
        \lexists[y][\psi(y)]}, just=\TAss, checked
      [\sFmla{\True}{\lnot\lexists[y][\psi(y)]}, just=\TAss, checked
        [\sFmla{\False}{\lexists[y][\psi(y)]},
          just={\TRule{\True}{\lnot}[3]},
          [\sFmla{\False}{\lforall[x][\varphi(x)]}, just={\TRule{\True}{\lif}[2]}]
          [\sFmla{\True}{\lexists[y][\psi(y)]}, just={\TRule{\True}{\lif}[2]}]
        ]
      ]
    ]
  ]
\end{oltableau}
दोन्ही नव्या चिन्हांकित सूत्रे ना आयगेन चराच्या अटी असलेले नियम आवश्यक
आहेत; म्हणून ते नियम यानंतर लावावेत:
\begin{oltableau}
  [\sFmla{\True}{\lforall[x][\varphi(x)]}, just=\TAss
    [\sFmla{\True}{\lforall[x][\varphi(x)] \lif
        \lexists[y][\psi(y)]}, just=\TAss, checked
      [\sFmla{\True}{\lnot\lexists[y][\psi(y)]}, just=\TAss, checked
        [\sFmla{\False}{\lexists[y][\psi(y)]},
          just={\TRule{\True}{\lnot}[3]}
          [\sFmla{\False}{\lforall[x][\varphi(x)]}, just={\TRule{\True}{\lif}[2]},checked
            [\sFmla{\False}{\varphi(b)}, just={\TRule{\False}{\lforall}[5]}]
          ]
          [\sFmla{\True}{\lexists[y][\psi(y)]}, just={\TRule{\True}{\lif}[2]},checked
            [\sFmla{\True}{\psi(c)}, just={\TRule{\True}{\lexists}[5]}]
          ]
        ]
      ]
    ]
  ]
\end{oltableau}
शाखा बंद करण्यासाठी ओळी $1$ आणि~$4$ वरील चिन्हांकित सूत्रे वापराव्या
लागतात. त्यांच्या संबंधित नियमांना (\TRule{\True}{\lforall} आणि
\TRule{\False}{\lexists}) आयगेन चराची अट नाही; म्हणून योग्य असलेली कोणतीही
बंद पदे निवडण्याचे स्वातंत्र्य आहे. येथे ती अनुक्रमे $b$ आणि~$c$ आहेत.
\begin{oltableau}
  [\sFmla{\True}{\lforall[x][\varphi(x)]}, just=\TAss
    [\sFmla{\True}{\lforall[x][\varphi(x)] \lif
        \lexists[y][\psi(y)]}, just=\TAss, checked
      [\sFmla{\True}{\lnot\lexists[y][\psi(y)]}, just=\TAss, checked
        [\sFmla{\False}{\lexists[y][\psi(y)]},
          just={\TRule{\True}{\lnot}[3]}
          [\sFmla{\False}{\lforall[x][\varphi(x)]},
            just={\TRule{\True}{\lif}[2]},checked
            [\sFmla{\False}{\varphi(b)},
              just={\TRule{\False}{\lforall}[5]}
              [\sFmla{\True}{\varphi(b)},
                just={\TRule{\True}{\lforall}[1]},close
              ]
            ]
          ]
          [\sFmla{\True}{\lexists[y][\psi(y)]},
            just={\TRule{\True}{\lif}[2]},checked
            [\sFmla{\True}{\psi(c)},
              just={\TRule{\True}{\lexists}[5]}
              [\sFmla{\False}{\psi(c)},
                just={\TRule{\False}{\lexists}[4]},close
              ]
            ]
          ]
        ]
      ]
    ]
  ]
\end{oltableau}
\end{ex}

\begin{prob}
पुढील सूत्रांचे बंद टॅब्लो द्या:
\begin{enumerate}
\item $\sFmla{\False}{(\lforall[x][\varphi(x)] \land \lforall[y][\psi(y)])
\lif \lforall[z][(\varphi(z) \land \psi(z))]}$.
\item $\sFmla{\False}{(\lexists[x][\varphi(x)] \lor \lexists[y][\psi(y)])
\lif \lexists[z][(\varphi(z) \lor \psi(z))]}$.
\item $\sFmla{\True}{\lforall[x][(\varphi(x) \lif \psi)]},
\sFmla{\False}{\lexists[y][\varphi(y)] \lif \psi}$.
\item $\sFmla{\True}{\lforall[x][\lnot \varphi(x)]},
\sFmla{\False}{\lnot\lexists[x][\varphi(x)]}$.
\item $\sFmla{\False}{\lnot\lexists[x][\varphi(x)] \lif \lforall[x][\lnot
\varphi(x)]}$.
\item $\sFmla{\False}{\lnot\lexists[x][\lforall[y][((\varphi(x,y) \lif
\lnot \varphi(y,y)) \land (\lnot \varphi(y,y) \lif \varphi(x,y)))]]}$.
\end{enumerate}
\end{prob}

\begin{prob}
पुढील सूत्रांचे बंद टॅब्लो द्या:
\begin{enumerate}
\item $\sFmla{\False}{\lnot\lforall[x][\varphi(x)] \lif \lexists[x][\lnot\varphi(x)]}$.
\item $\sFmla{\True}{(\lforall[x][\varphi(x)] \lif \psi)}, \sFmla{\False}{\lexists[y][(\varphi(y) \lif \psi)]}$.
\item $\sFmla{\False}{\lexists[x][(\varphi(x) \lif \lforall[y][\varphi(y)])]}$.
\end{enumerate}
\end{prob}












\section{सिद्धता-उपपत्तीय संकल्पना}\label{fol:tab:ptn:sec}

\begin{editorial}
या विभागात टॅब्लोसाठी सिद्धतायोग्यता संबंध आणि सुसंगतता यांच्या व्याख्या
एकत्रित केल्या आहेत.
\end{editorial}

\begin{explain}
जशा आपण अनेक महत्त्वाच्या चिन्हार्थविषयक संकल्पना (वैधता, तार्किक निष्पन्नता,
पूर्ततायोग्यता) परिभाषित केल्या आहेत, तशाच त्यांना अनुरूप
\emph{सिद्धता-उपपत्तीय संकल्पना} आता परिभाषित करू. संरचना मध्ये
वाक्यांची पूर्ति होते की नाही याचा आधार घेऊन या संकल्पना परिभाषित
केल्या जात नाहीत; त्याऐवजी ठरावीक बंद टॅब्लोx च्या अस्तित्वाचा आधार
घेतला जातो. या दोन प्रकारच्या संकल्पना एकमेकांशी जुळतात, हा एक महत्त्वाचा
शोध होता. त्या जुळतात, हाच \emph{निर्दोषता प्रमेय} आणि
\emph{संपूर्णता प्रमेय} यांचा आशय आहे.
\end{explain}


\begin{defn}[प्रमेये]
एखादे वाक्य~$\varphi$ हे \emph{प्रमेय} असते, जर
$\sFmla{\False}{\varphi}$ साठी बंद टॅब्लो असेल. $\varphi$ हे प्रमेय असेल
तर आपण $\Proves \varphi$ लिहितो आणि ते प्रमेय नसेल तर $\Proves/ \varphi$ लिहितो.
\end{defn}

\begin{defn}[निष्पन्नता]
वाक्य $\varphi$ हे वाक्यांच्या संच~$\Gamma$ \emph{पासून
निष्पन्न करता येण्याजोगे} असते, म्हणजे $\Gamma \Proves \varphi$, तेव्हा आणि केवळ तेव्हाच,
जेव्हा एखादा सांत संच $\{\psi_1, \dots, \psi_n\} \subseteq \Gamma$ असा असतो
आणि पुढील संचासाठी बंद टॅब्लो असतो:
\[\{
  \sFmla{\False}{\varphi},
  \sFmla{\True}{\psi_1}, \dots,
  \sFmla{\True}{\psi_n}
  \}.
\]
$\varphi$ हे $\Gamma$ पासून निष्पन्न करता येण्याजोगे नसेल, तर आपण $\Gamma \Proves/
\varphi$ लिहितो.
\end{defn}

\begin{defn}[सुसंगतता]
वाक्यांचा संच~$\Gamma$ हा \emph{विसंगत} असतो तेव्हा आणि केवळ तेव्हाच,
जेव्हा एखादा सांत संच $\{\psi_1, \dots, \psi_n\} \subseteq \Gamma$ असा असतो
आणि पुढील संचासाठी बंद टॅब्लो असतो:
\[\{
  \sFmla{\True}{\psi_1}, \dots,
  \sFmla{\True}{\psi_n}
  \}.
\]
$\Gamma$ विसंगत नसेल, तर त्याला \emph{सुसंगत} म्हणतो.
\end{defn}

\begin{prop}[परावर्तिता]
\label{fol:tab:ptn:prop:reflexivity}
$\varphi \in \Gamma$ असेल, तर $\Gamma \Proves \varphi$.
\end{prop}

\begin{proof}
$\varphi \in \Gamma$ असेल, तर $\{\varphi\}$ हा $\Gamma$ चा सांत उपसंच आहे आणि पुढील टॅब्लो
\begin{oltableau}
  [\sFmla{\False}{\varphi}, just = \TAss
    [\sFmla{\True}{\varphi}, just = \TAss,close]
  ]
\end{oltableau}
बंद आहे.
\end{proof}


\begin{prop}[एकस्वनिकता]
\label{fol:tab:ptn:prop:monotonicity}
$\Gamma \subseteq \Delta$ आणि $\Gamma \Proves \varphi$ असेल, तर $\Delta
\Proves \varphi$.
\end{prop}

\begin{proof}
$\Gamma$ चा कोणताही सांत उपसंच हा $\Delta$ चासुद्धा सांत उपसंच असतो.
\end{proof}

\begin{prop}[संक्रमकता]
\label{fol:tab:ptn:prop:transitivity}
$\Gamma \Proves \varphi$ आणि $\{\varphi\} \cup
\Delta \Proves \psi$ असेल, तर $\Gamma \cup \Delta \Proves \psi$.
\end{prop}

\begin{proof}
$\{\varphi\} \cup \Delta \Proves \psi$ असेल, तर $\Delta_0 =
\{\chi_1, \dots, \chi_n\} \subseteq \Delta$ असा सांत उपसंच असतो आणि पुढील संचासाठी
\begin{align*}
\{\sFmla{\False}{\psi}, & \sFmla{\True}{\varphi}, \sFmla{\True}{\chi_1},
\dots, \sFmla{\True}{\chi_n}\}
\intertext{बंद टॅब्लो असतो. जर $\Gamma \Proves \varphi$ असेल, तर
  $\{\theta_1, \dots, \theta_m\} \subseteq \Gamma$ असा संच असतो आणि पुढील संचासाठीही}
\{\sFmla{\False}{\varphi}, & \sFmla{\True}{\theta_1},
\dots, \sFmla{\True}{\theta_m}\}
\end{align*}
बंद टॅब्लो असतो.

आता पुढील गृहीतके असलेला टॅब्लो विचारात घ्या:
\[\sFmla{\False}{\psi},
\sFmla{\True}{\chi_1}, \dots, \sFmla{\True}{\chi_n},
\sFmla{\True}{\theta_1}, \dots, \sFmla{\True}{\theta_m}.
\]
त्यात~$\varphi$ वर \Cut{} नियम लावा. यामुळे दोन शाखा तयार होतात: एका शाखेत
$\sFmla{\True}{\varphi}$, तर दुसऱ्या शाखेत $\sFmla{\False}{\varphi}$ असते. त्यामुळे
पहिल्या शाखेत पुढील सर्व गृहीतके उपलब्ध असतात:
\[\{\sFmla{\False}{\psi}, \sFmla{\True}{\varphi},
\sFmla{\True}{\chi_1}, \dots, \sFmla{\True}{\chi_n}\}
\]
या गृहीतकांसाठी बंद टॅब्लो असल्यामुळे तो त्या शाखेला जोडता येतो;
$\sFmla{\True}{\varphi}$ मधून जाणारी प्रत्येक शाखा बंद होते. दुसऱ्या शाखेत
पुढील सर्व गृहीतके उपलब्ध असतात:
\[\{\sFmla{\False}{\varphi}, \sFmla{\True}{\theta_1}, \dots,
\sFmla{\True}{\theta_m}\}
\]
म्हणून दुसरी शाखासुद्धा पूर्ण करून बंद टॅब्लो मिळवता येतो. यावरून
$\Gamma \cup \Delta \Proves \psi$ हे दिसते.
\end{proof}

विशेषतः, याचा अर्थ असा की $\Gamma \Proves \varphi$ आणि $\varphi
\Proves \psi$ असेल, तर $\Gamma \Proves \psi$. तसेच $\varphi_1,
\dots, \varphi_n \Proves \psi$ असेल आणि प्रत्येक~$i$ साठी $\Gamma \Proves \varphi_i$
असेल, तर $\Gamma \Proves \psi$, हेसुद्धा यावरून मिळते.

\begin{prop}
\label{fol:tab:ptn:prop:incons}
$\Gamma$ विसंगत असतो तेव्हा आणि केवळ तेव्हाच, जेव्हा प्रत्येक
  वाक्य~$\varphi$ साठी $\Gamma \Proves \varphi$.
\end{prop}

\begin{proof}
सराव.
\end{proof}


\begin{prob}
\ref{fol:tab:ptn:prop:incons} सिद्ध करा.
\end{prob}




\begin{prop}[संहतता]
  \label{fol:tab:ptn:prop:proves-compact}
  \begin{enumerate}
  \item $\Gamma \Proves \varphi$ असेल, तर एखादा सांत उपसंच $\Gamma_0
    \subseteq \Gamma$ असा आहे की $\Gamma_0 \Proves \varphi$.
  \item $\Gamma$ चा प्रत्येक सांत उपसंच सुसंगत असेल, तर $\Gamma$ सुसंगत
    असतो.
  \end{enumerate}
\end{prop}

\begin{proof}
  \begin{enumerate}
    \item $\Gamma \Proves \varphi$ असेल, तर $\Gamma_0 =
      \{\psi_1, \dots, \psi_n\}$ असा त्याचा सांत उपसंच आणि पुढील संचासाठी बंद टॅब्लो असतो:
      \[      \{\sFmla{\False}{\varphi}, \sFmla{\True}{\psi_1}, \dots, \sFmla{\True}{\psi_n}\}
      \]
      हा टॅब्लो $\Gamma_0 \Proves \varphi$ हेसुद्धा दाखवतो.
    \item $\Gamma$ विसंगत असेल, तर
      $\Gamma_0 = \{\psi_1, \dots, \psi_n\}$ असा त्याचा एखादा सांत उपसंच असून पुढील
      संचाचा बंद टॅब्लो असतो:
      \[      \{\sFmla{\True}{\psi_1}, \dots, \sFmla{\True}{\psi_n}\}
      \]
      हा बंद टॅब्लो $\Gamma_0$ विसंगत असल्याचे दाखवतो.
  \end{enumerate}
\end{proof}












\section{निष्पन्नता आणि सुसंगतता}\label{fol:tab:prv:sec}

आता आपण निष्पन्नता संबंधाचे अनेक गुणधर्म सिद्ध करू. ते स्वतंत्रपणेही
महत्त्वाचे आहेत, पण संपूर्णता प्रमेयाच्या सिद्धतेत प्रत्येकाची भूमिका असेल.

\begin{prop}\label{fol:tab:prv:prop:provability-contr}
  जर $\Gamma \Proves \varphi$ आणि $\Gamma \cup \{\varphi\}$ विसंगत असेल, तर
  $\Gamma$ विसंगत आहे.
\end{prop}

\begin{proof}
$\Gamma_0 = \{\psi_1, \dots, \psi_n\}$ आणि $\Gamma_1
  =\{\chi_1, \dots, \chi_m\} \subseteq \Gamma$ असे सांत उपसंच असतात की
  \begin{align*}
    \{\sFmla{\False}{\varphi}, &
    \sFmla{\True}{\psi_1}, \dots, \sFmla{\True}{\psi_n}\} \\
    \{\sFmla{\True}{\varphi}, &
    \sFmla{\True}{\chi_1}, \dots, \sFmla{\True}{\chi_m}\}
  \end{align*}
  यांसाठी बंद टॅब्लो असतात. $\varphi$ वर \Cut{} नियम वापरून हे दोन्ही
  एकाच बंद टॅब्लोमध्ये एकत्र करता येतात; त्यावरून $\Gamma_0
  \cup \Gamma_1$ विसंगत असल्याचे दिसते. $\Gamma_0
  \subseteq \Gamma$ आणि $\Gamma_1 \subseteq \Gamma$ असल्यामुळे $\Gamma_0 \cup
  \Gamma_1 \subseteq \Gamma$; म्हणून $\Gamma$ विसंगत आहे.
\end{proof}

\begin{prop}
\label{fol:tab:prv:prop:prov-incons}
$\Gamma \Proves \varphi$ तेव्हा आणि केवळ तेव्हाच, जेव्हा $\Gamma \cup \{\lnot \varphi\}$
विसंगत असते.
\end{prop}

\begin{proof}
प्रथम $\Gamma \Proves \varphi$ असे समजा; म्हणजे पुढील संचासाठी बंद
टॅब्लो असतो:
\[\{\sFmla{\False}{\varphi},
\sFmla{\True}{\psi_1}, \dots, \sFmla{\True}{\psi_n}\}
\]
$\TRule{\True}{\lnot}$ नियम वापरून त्याचे पुढील संचासाठीच्या बंद
टॅब्लोमध्ये रूपांतर करता येते:
\[\{\sFmla{\True}{\lnot \varphi},
\sFmla{\True}{\psi_1}, \dots, \sFmla{\True}{\psi_n}\}.
\]

दुसरीकडे, दुसऱ्या संचासाठी बंद टॅब्लो असेल, तर त्यातील पहिले
गृहीतक~$\sFmla{\True}{\lnot \varphi}$ आणि त्या गृहीतकाला
\TRule{\True}{\lnot} लावून मिळणारी सर्व सूत्रे काढून, तसेच
$\sFmla{\False}{\varphi}$ हे गृहीतक जोडून, त्याचे पहिल्या संचासाठीच्या बंद
टॅब्लोमध्ये रूपांतर करता येते. कारण एखादी शाखा पूर्वी
$\sFmla{\True}{\lnot \varphi}$ ला \TRule{\True}{\lnot} लावून मिळणारा
निष्कर्ष, म्हणजे $\sFmla{\False}{\varphi}$, असल्यामुळे बंद झाली असेल, तर नव्या
टॅब्लो मधील संबंधित शाखाही बंद असते. जुन्या टॅब्लोतील एखादी शाखा
$\sFmla{\True}{\lnot \varphi}$ हे गृहीतक आणि $\sFmla{\False}{\lnot \varphi}$
ही दोन्ही सूत्रे असल्यामुळे बंद झाली असेल, तर $\sFmla{\False}{\lnot \varphi}$
ला $\TRule{\False}{\lnot}$ लावून $\sFmla{\True}{\varphi}$ मिळवता येते आणि
ती शाखा बंद करता येते. आपण $\sFmla{\False}{\varphi}$ हे गृहीतक जोडलेले
असल्यामुळे ही शाखा बंद होते.
\end{proof}

\begin{prob}
$\Gamma \Proves \lnot \varphi$ तेव्हा आणि केवळ तेव्हाच, जेव्हा $\Gamma \cup \{\varphi\}$
विसंगत असते, हे सिद्ध करा.
\end{prob}

\begin{prop}\label{fol:tab:prv:prop:explicit-inc}
  जर $\Gamma \Proves \varphi$ आणि $\lnot \varphi \in \Gamma$ असेल, तर $\Gamma$
  विसंगत आहे.
\end{prop}

\begin{proof}
  समजा $\Gamma \Proves \varphi$ आणि $\lnot \varphi \in \Gamma$. मग
  $\psi_1$, \dots, $\psi_n \in \Gamma$ अशी वाक्ये असतात की
  \[  \{\sFmla{\False}{\varphi}, \sFmla{\True}{\psi_1}, \dots, \sFmla{\True}{\psi_n}\}
  \]
  यासाठी बंद टॅब्लो असतो. \sFmla{\False}{\varphi} हे गृहीतक
  \sFmla{\True}{\lnot \varphi} ने बदला आणि \sFmla{\True}{\lnot \varphi} ला
  \TRule{\True}{\lnot} लावून मिळणारा निष्कर्ष गृहीतकांनंतर घाला. जुन्या
  टॅब्लोमध्ये ओळ~$1$ च्या आधारे समर्थित कोणतेही वाक्य आता
  ओळ~$n+2$ च्या आधारे समर्थित होते. म्हणून जुना टॅब्लो बंद असेल, तर
  नवा टॅब्लोसुद्धा बंद असतो. सर्व गृहीतके~$\Gamma$ मध्ये असल्यामुळे हा टॅब्लो
  $\Gamma$ विसंगत असल्याचे दाखवतो.
\end{proof}

\begin{prop}\label{fol:tab:prv:prop:provability-exhaustive}
  $\Gamma \cup \{\varphi\}$ आणि $\Gamma \cup \{\lnot \varphi\}$ हे दोन्ही संच
  विसंगत असतील, तर $\Gamma$ विसंगत आहे.
\end{prop}

\begin{proof}
  $\psi_1$, \dots, $\psi_n \in \Gamma$ आणि $\chi_1$, \dots,
  $\chi_m \in \Gamma$ अशी वाक्ये असतील आणि
  \begin{align*}
    \{\sFmla{\True}{\varphi}, &
    \sFmla{\True}{\psi_1}, \dots, \sFmla{\True}{\psi_n}\} \text{ आणि}\\
    \{\sFmla{\True}{\lnot \varphi}, &
    \sFmla{\True}{\chi_1}, \dots, \sFmla{\True}{\chi_m}\}
  \end{align*}
  या दोन्हींसाठी बंद टॅब्लो असतील, तर
  $\sFmla{\True}{\psi_1}$, \dots, $\sFmla{\True}{\psi_n}$ आणि
  $\sFmla{\True}{\chi_1}$, \dots, $\sFmla{\True}{\chi_m}$ ही गृहीतके घेऊन,
  त्यानंतर \Cut{} नियम लावून, $\Gamma$ विसंगत असल्याचे दाखवणारा एकच
  संयुक्त टॅब्लो तयार करता येतो. त्यामुळे दोन शाखा मिळतात: एक
  $\sFmla{\True}{\varphi}$ ने, तर दुसरी $\sFmla{\False}{\varphi}$ ने सुरू होते.

  डाव्या बाजूस पहिल्या टॅब्लो मधील त्याच्या गृहीतकांखालचा भाग जोडा.
  येथे प्रत्येक नियमाचा उपयोग अजूनही योग्य आहे, कारण
  $\sFmla{\True}{\varphi}$ सहित पहिल्या टॅब्लो मधील प्रत्येक गृहीतक
  उपलब्ध आहे. त्यामुळे $\sFmla{\True}{\varphi}$ च्या खालील प्रत्येक शाखा बंद होते.

  उजव्या बाजूस दुसऱ्या टॅब्लो मधील त्याच्या गृहीतकांखालचा भाग जोडा;
  मात्र $\sFmla{\True}{\lnot \varphi}$ ला~$\TRule{\True}{\lnot}$ लावून मिळणारे
  सर्व परिणाम काढून टाका. $\sFmla{\True}{\lnot \varphi}$ ला
  $\TRule{\True}{\lnot}$ लावल्याचा निष्कर्ष $\sFmla{\False}{\varphi}$ हा आहे;
  तरीही तो उपलब्ध असतो, कारण तो संयुक्त टॅब्लोच्या उजव्या बाजूवरील
  \Cut{} नियमाचा निष्कर्ष आहे.

  दुसऱ्या टॅब्लोतील एखादी शाखा $\sFmla{\True}{\lnot \varphi}$ हे गृहीतक
  (जे आता संयुक्त टॅब्लोमध्ये गृहीतक म्हणून येत नाही) आणि
  $\sFmla{\False}{\lnot \varphi}$ ही दोन्ही सूत्रे असल्यामुळे बंद झाली असेल,
  तर $\sFmla{\False}{\lnot \varphi}$ ला $\TRule{\False}{\lnot}$ लावून
  $\sFmla{\True}{\varphi}$ मिळवता येते. मग संयुक्त टॅब्लो मधील संबंधित
  शाखासुद्धा बंद होते, कारण तिच्यात \Cut{} नियमाचा उजवीकडील निष्कर्ष
  $\sFmla{\False}{\varphi}$ असतो. दुसऱ्या टॅब्लो मधील एखादी शाखा अन्य
  कोणत्याही कारणाने बंद झाली असेल, तर संयुक्त टॅब्लो मधील संबंधित
  शाखासुद्धा बंद होते; कारण जुन्या दुसऱ्या टॅब्लो मधील त्या शाखेवर
  येणारी $\sFmla{\True}{\lnot \varphi}$ व्यतिरिक्त इतर सर्व चिन्हांकित सूत्रे
  संयुक्त टॅब्लो मधील संबंधित शाखेवरही येतात.
  \end{proof}















\section{निष्पन्नता आणि विधानीय संयोजक}\label{fol:tab:ppr:sec}

\begin{explain}
  टॅब्लोचा निष्पन्नता संबंध~$\Proves$ हा विधानीय संयोजकांशी संबंधित
  काही मूलभूत तथ्ये सिद्ध करण्यासाठी पुरेसा सामर्थ्यवान आहे; उदाहरणार्थ,
  $\varphi \land \psi \Proves \varphi$ आणि $\varphi, \varphi \lif \psi \Proves \psi$
  (विध्यात्मक अनुमान). ही तथ्ये संपूर्णता प्रमेयाच्या सिद्धतेसाठी आवश्यक आहेत.
\end{explain}

\begin{prop}\label{fol:tab:ppr:prop:provability-land}
  \begin{enumerate}
  \item \label{fol:tab:ppr:prop:provability-land-left} $\varphi \land \psi \Proves
    \varphi$ आणि $\varphi \land \psi \Proves \psi$ ही दोन्ही तथ्ये सत्य आहेत.
  \item \label{fol:tab:ppr:prop:provability-land-right} $\varphi, \psi \Proves \varphi \land
    \psi$.
  \end{enumerate}
\end{prop}

\begin{proof}
  \begin{enumerate}
  \item $\{\sFmla{\False}{\varphi}, \sFmla{\True}{\varphi \land \psi}\}$ आणि
    $\{\sFmla{\False}{\psi}, \sFmla{\True}{\varphi \land \psi}\}$ या दोन्हींसाठी
    बंद टॅब्लो आहेत:
    \begin{oltableau}{}
      [\sFmla{\False}{\varphi}, just=\TAss
        [\sFmla{\True}{\varphi \land \psi}, just=\TAss
          [\sFmla{\True}{\varphi},just={\TRule{\True}{\land}[2]}
            [\sFmla{\True}{\psi},just={\TRule{\True}{\land}[2]}, close
            ]
          ]
        ]
      ]
    \end{oltableau}
    \begin{oltableau}{}
      [\sFmla{\False}{\psi}, just=\TAss
        [\sFmla{\True}{\varphi \land \psi}, just=\TAss
          [\sFmla{\True}{\varphi},just={\TRule{\True}{\land}[2]}
            [\sFmla{\True}{\psi},just={\TRule{\True}{\land}[2]}, close
            ]
          ]
        ]
      ]
    \end{oltableau}
    \item $\{\sFmla{\True}{\varphi}, \sFmla{\True}{\psi},
      \sFmla{\False}{\varphi \land \psi}\}$ साठीचा बंद टॅब्लो पुढीलप्रमाणे आहे:
      \begin{oltableau}
        [\sFmla{\False}{\varphi \land \psi}, just = \TAss
          [\sFmla{\True}{\varphi}, just = \TAss
            [\sFmla{\True}{\psi}, just=\TAss
              [\sFmla{\False}{\varphi}, just = {\TRule{\False}{\land}[1]}, close]
              [\sFmla{\False}{\psi}, just = {\TRule{\False}{\land}[1]}, close]
            ]
          ]
        ]
      \end{oltableau}
  \end{enumerate}
\end{proof}

\begin{prop}\label{fol:tab:ppr:prop:provability-lor}
  \begin{enumerate}
  \item $\{\varphi \lor \psi, \lnot \varphi, \lnot \psi\}$ विसंगत आहे.
  \item $\varphi \Proves \varphi \lor \psi$ आणि $\psi \Proves \varphi \lor \psi$ ही दोन्ही
    तथ्ये सत्य आहेत.
  \end{enumerate}
\end{prop}

\begin{proof}
  \begin{enumerate}
  \item $\{\sFmla{\True}{\varphi \lor \psi}, \sFmla{\True}{\lnot \varphi},
    \sFmla{\True}{\lnot \psi}\}$ साठी बंद टॅब्लो देऊ:
    \begin{oltableau}
      [\sFmla{\True}{\varphi \lor \psi}, just = \TAss
        [\sFmla{\True}{\lnot \varphi}, just = \TAss
          [\sFmla{\True}{\lnot \psi}, just = \TAss
            [\sFmla{\False}{\varphi}, just = {\TRule{\True}{\lnot}[2]}
              [\sFmla{\False}{\psi}, just = {\TRule{\True}{\lnot}[3]}
                [\sFmla{\True}{\varphi}, just = {\TRule{\True}{\lor}[1]}, close]
                [\sFmla{\True}{\psi}, just = {\TRule{\True}{\lor}[1]}, close]
              ]
            ]
          ]
        ]
      ]
    \end{oltableau}
  \item $\{\sFmla{\False}{\varphi \lor \psi}, \sFmla{\True}{\varphi}\}$ आणि
    $\{\sFmla{\False}{\varphi \lor \psi}, \sFmla{\True}{\psi}\}$ या दोन्हींसाठी
    बंद टॅब्लो आहेत:
    \begin{oltableau}{}
      [\sFmla{\False}{\varphi \lor \psi}, just=\TAss
        [\sFmla{\True}{\varphi}, just=\TAss
          [\sFmla{\False}{\varphi},just={\TRule{\False}{\lor}[1]}
            [\sFmla{\False}{\psi},just={\TRule{\False}{\lor}[1]}, close
            ]
          ]
        ]
      ]
    \end{oltableau}
    \begin{oltableau}{}
      [\sFmla{\False}{\varphi \lor \psi}, just=\TAss
        [\sFmla{\True}{\psi}, just=\TAss
          [\sFmla{\False}{\varphi},just={\TRule{\False}{\lor}[1]}
            [\sFmla{\False}{\psi},just={\TRule{\False}{\lor}[1]}, close
            ]
          ]
        ]
      ]
    \end{oltableau}
  \end{enumerate}
\end{proof}

\begin{prop}\label{fol:tab:ppr:prop:provability-lif}
  \begin{enumerate}
  \item \label{fol:tab:ppr:prop:provability-lif-left}  $\varphi, \varphi \lif \psi \Proves \psi$.
  \item \label{fol:tab:ppr:prop:provability-lif-right}
    $\lnot \varphi \Proves \varphi \lif \psi$ आणि $\psi \Proves \varphi \lif \psi$ ही दोन्ही
    तथ्ये सत्य आहेत.
  \end{enumerate}
\end{prop}

\begin{proof}
  \begin{enumerate}
    \item $\{\sFmla{\False}{\psi}, \sFmla{\True}{\varphi \lif \psi},
      \sFmla{\True}{\varphi}\}$ साठी बंद टॅब्लो आहे:
      \begin{oltableau}
        [\sFmla{\False}{\psi}, just=\TAss
          [\sFmla{\True}{\varphi \lif \psi}, just=\TAss
            [\sFmla{\True}{\varphi}, just=\TAss
              [\sFmla{\False}{\varphi}, just = {\TRule{\True}{\lif}[2]}, close]
              [\sFmla{\True}{\psi}, just = {\TRule{\True}{\lif}[2]}, close]
            ]
          ]
        ]
      \end{oltableau}
    \item $\{\sFmla{\False}{\varphi \lif \psi}, \sFmla{\True}{\lnot \varphi}\}$ आणि
      $\{\sFmla{\False}{\varphi \lif \psi}, \sFmla{\True}{\psi}\}$ या दोन्हींसाठी
      बंद टॅब्लो आहेत:
      \begin{oltableau}
        [\sFmla{\False}{\varphi \lif \psi}, just = \TAss
          [\sFmla{\True}{\lnot \varphi}, just = \TAss
            [\sFmla{\True}{\varphi}, just = {\TRule{\False}{\lif}[1]}
              [\sFmla{\False}{\psi}, just = {\TRule{\False}{\lif}[1]}
                [\sFmla{\False}{\varphi}, just = {\TRule{\True}{\lnot}[2]}, close]
              ]
            ]
          ]
        ]
      \end{oltableau}
      \begin{oltableau}
        [\sFmla{\False}{\varphi \lif \psi}, just = \TAss
          [\sFmla{\True}{\psi}, just = \TAss
            [\sFmla{\True}{\varphi}, just = {\TRule{\False}{\lif}[1]}
              [\sFmla{\False}{\psi}, just = {\TRule{\False}{\lif}[1]},close
              ]
            ]
          ]
        ]
      \end{oltableau}
  \end{enumerate}
\end{proof}











\section{निष्पन्नता आणि संख्यापक}\label{fol:tab:qpr:sec}

\begin{explain}
  संपूर्णता प्रमेयासाठी या विभागात स्थापित केलेली~$\Proves$ संबंधी तथ्ये
  टॅब्लोच्या नियमांनी निष्पन्न होतात, हेही आवश्यक आहे.
\end{explain}

\begin{thm}
\label{fol:tab:qpr:thm:strong-generalization} जर $c$ हे $\Gamma$ किंवा~$\varphi(x)$ मध्ये
न येणारे अचर असेल आणि $\Gamma \Proves \varphi(c)$, तर $\Gamma
\Proves \lforall[x][\varphi(x)]$.
\end{thm}

\begin{proof}
$\Gamma \Proves \varphi(c)$ असे समजा; म्हणजे $\psi_1$, \dots, $\psi_n
\in \Gamma$ अशी वाक्ये आणि पुढील संचासाठी बंद टॅब्लो असतो:
\begin{align*}
\{ \sFmla{\False}{\varphi(c)}, &
\sFmla{\True}{\psi_1}, \dots, \sFmla{\True}{\psi_n} \}.
\intertext{पुढील संचासाठीही बंद टॅब्लो असतो, हे दाखवायचे आहे:}
\{ \sFmla{\False}{\lforall[x][\varphi(x)]}, &
\sFmla{\True}{\psi_1}, \dots, \sFmla{\True}{\psi_n} \}.
\end{align*}
बंद टॅब्लो घेऊन त्यातील पहिले गृहीतक
$\sFmla{\False}{\lforall[x][\varphi(x)]}$ ने बदला आणि गृहीतकांनंतर
$\sFmla{\False}{\varphi(c)}$ घाला.
\begin{center}
\begin{tableau}{not line numbering}
  [\sFmla{\False}{\varphi(c)}
    [\sFmla{\True}{\psi_1}
      [\vdots
      [\sFmla{\True}{\psi_n},
        [][][]]]]]
\end{tableau}\qquad
\begin{tableau}{not line numbering}
  [\sFmla{\False}{\lforall[x][\varphi(x)]}
    [\sFmla{\True}{\psi_1}
      [\vdots
        [\sFmla{\True}{\psi_n},
          [\sFmla{\False}{\varphi(c)}
        [][][]]]]]]
\end{tableau}
\end{center}
हा टॅब्लो अजूनही बंद आहे, कारण पूर्वी गृहीतके म्हणून उपलब्ध असलेली सर्व
वाक्ये नव्या टॅब्लोच्या वरच्या भागात अजूनही उपलब्ध आहेत. घातलेली
ओळ ही $\TRule{\False}{\lforall}$ च्या योग्य उपयोगाचा परिणाम आहे, कारण
स्थिरांक~$c$ हे $\psi_1$, \dots, $\psi_n$ किंवा
$\lforall[x][\varphi(x)]$ मध्ये येत नाही; म्हणजे ते नव्या टॅब्लोमध्ये घातलेल्या
ओळीच्या वर येत नाही.
\end{proof}

\begin{prop}
\label{fol:tab:qpr:prop:provability-quantifiers}
\begin{enumerate}
\item $\varphi(t) \Proves \lexists[x][\varphi(x)]$.
\item $\lforall[x][\varphi(x)] \Proves \varphi(t)$.
\end{enumerate}
\end{prop}

\begin{proof}
\begin{enumerate}
\item
  $\sFmla{\False}{\lexists[x][\varphi(x)]}, \sFmla{\True}{\varphi(t)}$ साठीचा बंद
  टॅब्लो पुढीलप्रमाणे आहे:
  \begin{oltableau}
    [\sFmla{\False}{\lexists[x][\varphi(x)]}, just = \TAss
      [\sFmla{\True}{\varphi(t)}, just = \TAss
        [\sFmla{\False}{\varphi(t)}, just = {\TRule{\False}{\lexists}[1]}, close]]]
    \end{oltableau}
\item
  $\sFmla{\False}{\varphi(t)}, \sFmla{\True}{\lforall[x][\varphi(x)]}, $
  साठीचा बंद टॅब्लो पुढीलप्रमाणे आहे:
  \begin{oltableau}
    [\sFmla{\False}{\varphi(t)}, just = \TAss
      [\sFmla{\True}{\lforall[x][\varphi(x)]}, just = \TAss
        [\sFmla{\True}{\varphi(t)}, just = {\TRule{\True}{\lforall}[2]}, close]]]
    \end{oltableau}
\end{enumerate}
\end{proof}












\section{निर्दोषता}\label{fol:tab:sou:sec}

\begin{explain}
टॅब्लोसारखी निष्पत्ती-पद्धत \emph{निर्दोष} असते, जर ती प्रत्यक्षात
धारण न करणाऱ्या गोष्टी निष्पन्न करू शकत नसेल. त्यामुळे निर्दोषता हा
निष्पत्ती-पद्धतींसाठी हमी दिलेल्या सुरक्षिततेचा एक प्रकार आहे. कोणता
सिद्धता-उपपत्तीय गुणधर्म विचारात आहे यावर अवलंबून, उदाहरणार्थ, आपल्याला
पुढील गोष्टी हव्या असतात:
\begin{enumerate}
\item प्रत्येक निष्पन्न करता येण्याजोगे~$\varphi$ हे वैध असते;
\item एखादे वाक्य इतर काही वाक्यांपासून निष्पन्न करता येण्याजोगे असेल, तर ते
  त्यांचा तार्किक परिणामही असते;
\item वाक्यांचा संच विसंगत असेल, तर तो अपूर्ततायोग्य असतो.
\end{enumerate}
हे निष्पत्ती-पद्धतीचे महत्त्वाचे गुणधर्म आहेत. यांपैकी एखादा गुणधर्म
खरा नसेल, तर निष्पत्ती-पद्धतीत दोष आहे---ती खूप जास्त गोष्टी
निष्पन्न करेल. म्हणून निष्पत्ती-पद्धतीची निर्दोषता सिद्ध करणे अत्यंत
महत्त्वाचे आहे.

हे सर्व सिद्धता-उपपत्तीय गुणधर्म कोणत्या ना कोणत्या प्रकारच्या बंद
टॅब्लो द्वारे परिभाषित केलेले असल्यामुळे, वरील (1)--(3) सिद्ध करण्यासाठी
बंद टॅब्लोच्या चिन्हार्थविषयक गुणधर्मांबद्दल काहीतरी सिद्ध करावे लागते.
प्रथम एखाद्या चिन्हांकित सूत्राची संरचनेत पूर्ती होणे म्हणजे काय ते
परिभाषित करू; मग टॅब्लो बंद असेल, तर कोणतीही संरचना त्याची सर्व गृहीतके
पूर्ण करत नाही, हे दाखवू. त्यानंतर (1)--(3) या परिणामाचे उपपरिणाम म्हणून
मिळतील.
\end{explain}

\begin{defn}
संरचना~$\Struct{M}$
चिन्हांकित सूत्र $\sFmla{\True}{\varphi}$ \emph{पूर्ण करते} तेव्हा आणि केवळ तेव्हाच, जेव्हा
$\Sat{M}{\varphi}$; आणि ते
$\sFmla{\False}{\varphi}$ पूर्ण करते तेव्हा आणि केवळ तेव्हाच, जेव्हा
$\Sat/{M}{\varphi}$. $\Struct{M}$
हा चिन्हांकित सूत्रांचा संच~$\Gamma$ पूर्ण करतो तेव्हा आणि केवळ तेव्हाच,
जेव्हा तो प्रत्येक $\sFmla{S}{\varphi} \in \Gamma$ पूर्ण करतो. $\Gamma$ पूर्ण
करणारी एखादी संरचना असेल, तर तो
\emph{पूर्ततायोग्य} आणि अन्यथा \emph{अपूर्ततायोग्य} असतो.
\end{defn}

\begin{thm}[निर्दोषता]
  \label{fol:tab:sou:thm:tableau-soundness}
  $\Gamma$ साठी बंद टॅब्लो असेल, तर $\Gamma$ अपूर्ततायोग्य असतो.
\end{thm}

\begin{proof}
टॅब्लोच्या शाखेवरील चिन्हांकित सूत्रांचा संच पूर्ततायोग्य असेल, तर
आणि केवळ तेव्हाच त्या शाखेला पूर्ततायोग्य म्हणू; तसेच टॅब्लोमध्ये
किमान एक पूर्ततायोग्य शाखा असेल, तर त्याला पूर्ततायोग्य म्हणू.

पुढील गोष्ट दाखवू: एखाद्या पूर्ततायोग्य टॅब्लोचा अनुमानाच्या कोणत्याही
एका नियमाने विस्तार केल्यास मिळणारा टॅब्लो नेहमी पूर्ततायोग्य असतो.
यावरून प्रमेय सिद्ध होईल: $\Gamma$ मधील केवळ गृहीतके असलेल्या टॅब्लो ला
अनुमानाचे नियम लावून कोणताही बंद टॅब्लो मिळतो. म्हणून $\Gamma$
पूर्ततायोग्य असेल, तर त्याच्यासाठीचा कोणताही टॅब्लो पूर्ततायोग्य असेल.
पण बंद टॅब्लो स्पष्टपणे पूर्ततायोग्य नसतो: प्रत्येक शाखेत
$\sFmla{\True}{\varphi}$ आणि $\sFmla{\False}{\varphi}$ ही दोन्ही सूत्रे असतात, आणि
कोणतीही संरचना $\varphi$ ची पूर्ती आणि अपूर्ती या दोन्ही करू शकत नाही.

आपल्याकडे पूर्ततायोग्य टॅब्लो, म्हणजे किमान एक पूर्ततायोग्य शाखा असलेला
टॅब्लो, आहे असे समजा. अनुमानाचा नियम लावल्यावर एकतर शाखेत
चिन्हांकित सूत्रे जोडली जातात, किंवा शाखा दोन भागांत विभागली जाते. त्या
नियमाच्या उपयोगाने विस्तारली न गेलेली एखादी पूर्ततायोग्य शाखा टॅब्लोमध्ये
असेल, तर विस्तारित टॅब्लोमध्येही ती पूर्ततायोग्य शाखा राहते; म्हणून
विस्तारित टॅब्लोसुद्धा पूर्ततायोग्य असतो. त्यामुळे पूर्ततायोग्य शाखेलाच नियम
लावला जातो ते प्रकरण पाहणे पुरेसे आहे.

त्या शाखेवरील चिन्हांकित सूत्रांचा संच~$\Gamma$ असू द्या, आणि ज्या
चिन्हांकित सूत्राला नियम लावला जातो ते $\sFmla{S}{\varphi} \in \Gamma$ असू
द्या. नियम लावल्यावर शाखा विभागली जात नसेल, तर विस्तारित शाखा, म्हणजे
$\Gamma$ आणि नियमाचे निष्कर्ष मिळून बनलेला संच, अजूनही पूर्ततायोग्य आहे हे
दाखवावे लागेल. नियम लावल्यावर शाखा विभागली जात असेल, तर मिळणाऱ्या दोन
शाखांपैकी किमान एक पूर्ततायोग्य आहे हे दाखवावे लागेल.

प्रथम, शाखा न विभागणारी संभाव्य अनुमाने पाहू.
\begin{enumerate}
\item $\sFmla{\True}{\lnot \psi} \in \Gamma$ ला
  $\TRule{\True}{\lnot}$ लावून शाखेचा विस्तार केला आहे असे समजा. मग
  विस्तारित शाखेवरील चिन्हांकित सूत्रांचा संच $\Gamma \cup
  \{\sFmla{\False}{\psi}\}$ असतो. समजा
  $\Sat{M}{\Gamma}$. विशेषतः,
  $\Sat{M}{\lnot \psi}$. त्यामुळे
  $\Sat/{M}{\psi}$, म्हणजे
  $\Struct{M}$ हे
  $\sFmla{\False}{\psi}$ पूर्ण करते.
\item $\sFmla{\False}{\lnot \psi} \in \Gamma$ ला
  $\TRule{\False}{\lnot}$ लावून शाखेचा विस्तार केला आहे: सराव.
\item $\sFmla{\True}{\psi \land \chi} \in \Gamma$ ला
  $\TRule{\True}{\land}$ लावून शाखेचा विस्तार केला आहे. यामुळे शाखेवर
  $\sFmla{\True}{\psi}$ आणि $\sFmla{\True}{\chi}$ ही दोन नवी
  चिन्हांकित सूत्रे मिळतात. समजा
  $\Sat{M}{\Gamma}$; विशेषतः,
  $\Sat{M}{\psi \land \chi}$. मग
  $\Sat{M}{\psi}$ आणि
  $\Sat{M}{\chi}$. म्हणजे
  $\Struct{M}$ हे
  $\sFmla{\True}{\psi}$ आणि $\sFmla{\True}{\chi}$ ही दोन्ही सूत्रे पूर्ण करते.
\item $\sFmla{\False}{\psi \lor \chi} \in \Gamma$ ला
  $\TRule{\False}{\lor}$ लावून शाखेचा विस्तार केला आहे: सराव.
\item $\sFmla{\False}{\psi \lif \chi} \in \Gamma$ ला
  $\TRule{\False}{\lif}$ लावून शाखेचा विस्तार केला आहे. यामुळे शाखेवर
  $\sFmla{\True}{\psi}$ आणि $\sFmla{\False}{\chi}$ ही दोन नवी
  चिन्हांकित सूत्रे मिळतात. समजा
  $\Sat{M}{\Gamma}$; विशेषतः,
  $\Sat/{M}{\psi \lif \chi}$. मग
  $\Sat{M}{\psi}$ आणि
  $\Sat/{M}{\chi}$. म्हणजे
  $\Struct{M}$ हे
  $\sFmla{\True}{\psi}$ आणि $\sFmla{\False}{\chi}$ ही दोन्ही सूत्रे पूर्ण करते.


\item $\sFmla{\True}{\lforall[x][\varphi(x)]} \in \Gamma$ ला
  $\TRule{\True}{\lforall}$ लावून शाखेचा विस्तार केला आहे. यामुळे शाखेवर
  नवीन चिन्हांकित सूत्र~$\sFmla{\True}{\varphi(t)}$ मिळते.
  समजा $\Sat{M}{\Gamma}$; विशेषतः,
  $\Sat{M}{\lforall[x][\varphi(x)]}$.
  \ref{fol:syn:sem:prop:quant-terms} नुसार $\Sat{M}{\varphi(t)}$. परिणामी,
  $\Struct{M}$ हे $\sFmla{\True}{\varphi(t)}$ पूर्ण करते.

\item $\sFmla{\False}{\lforall[x][\varphi(x)]} \in \Gamma$ ला
  $\TRule{\False}{\lforall}$ लावून शाखेचा विस्तार केला आहे. यामुळे शाखेवर
  नवीन चिन्हांकित सूत्र~$\sFmla{\False}{\varphi(a)}$ मिळते, जेथे $a$ हे
  $\Gamma$ मध्ये न येणारे स्थिरांक आहे. $\Gamma$ पूर्ततायोग्य असल्यामुळे
  $\Sat{M}{\Gamma}$ अशी एखादी $\Struct{M}$ असते; विशेषतः,
  $\Sat/{M}{\lforall[x][\varphi(x)]}$. आता
  $\Gamma \cup \{\sFmla{\False}{\varphi(a)}\}$ पूर्ततायोग्य आहे हे दाखवायचे आहे.
  त्यासाठी पुढीलप्रमाणे योग्य~$\Struct{M'}$ परिभाषित करू.

  \ref{fol:syn:ass:prop:sat-quant} नुसार,
  $\Sat/{M}{\lforall[x][\varphi(x)]}$ तेव्हा आणि केवळ तेव्हाच, जेव्हा काही $s$
  साठी $\Sat/{M}{\varphi(x)}[s]$. आता $\Struct{M'}$ ही $\Struct{M}$ सारखीच
  असू द्या, फक्त $\Assign{a}{M'} = s(x)$. $a$ हे $\Gamma$ मध्ये येत
  नसल्यामुळे, \ref{fol:syn:ext:cor:extensionality-sent} नुसार कोणत्याही
  $\sFmla{\True}{\chi} \in \Gamma$ साठी $\Sat{M'}{\chi}$, आणि कोणत्याही
  $\sFmla{\False}{\chi} \in \Gamma$ साठी $\Sat/{M'}{\chi}$.

  \ref{fol:syn:ext:prop:extensionality} नुसार $\Sat/{M'}{\varphi(x)}[s]$.
  \ref{fol:syn:ext:prop:ext-formulas} नुसार $\Sat/{M'}{\varphi(a)}[s]$.
  $\varphi(a)$ हे वाक्य असल्यामुळे,
  \ref{fol:syn:ass:prop:sentence-sat-true} नुसार $\Sat/{M'}{\varphi(a)}$;
  म्हणजे $\Struct{M'}$ हे $\sFmla{\False}{\varphi(a)}$ पूर्ण करते.
\item $\sFmla{\True}{\lexists[x][\psi(x)]} \in \Gamma$ ला
  $\TRule{\True}{\lexists}$ लावून शाखेचा विस्तार केला आहे: सराव.
\item $\sFmla{\False}{\lexists[x][\psi(x)]} \in \Gamma$ ला
  $\TRule{\False}{\lexists}$ लावून शाखेचा विस्तार केला आहे: सराव.
\end{enumerate}
आता शाखा विभागणारी संभाव्य अनुमाने पाहू.
\begin{enumerate}
\item $\sFmla{\False}{\psi \land \chi} \in \Gamma$ ला
  $\TRule{\False}{\land}$ लावून शाखेचा विस्तार केला आहे. यामुळे दोन शाखा
  मिळतात: डावीकडील शाखा $\sFmla{\False}{\psi}$ मधून आणि उजवीकडील शाखा
  $\sFmla{\False}{\chi}$ मधून पुढे जाते. समजा
  $\Sat{M}{\Gamma}$; विशेषतः,
  $\Sat/{M}{\psi \land \chi}$. मग एकतर
  $\Sat/{M}{\psi}$ किंवा
  $\Sat/{M}{\chi}$. पहिल्या प्रकरणात
  $\Struct{M}$ हे
  $\sFmla{\False}{\psi}$ पूर्ण करते; म्हणजे
  $\Struct{M}$ डावीकडील शाखेवरील सूत्रे पूर्ण करते.
  दुसऱ्या प्रकरणात $\Struct{M}$ हे
  $\sFmla{\False}{\chi}$ पूर्ण करते; म्हणजे
  $\Struct{M}$ उजवीकडील शाखेवरील सूत्रे पूर्ण करते.
\item $\sFmla{\True}{\psi \lor \chi} \in \Gamma$ ला
  $\TRule{\True}{\lor}$ लावून शाखेचा विस्तार केला आहे: सराव.
\item $\sFmla{\True}{\psi \lif \chi} \in \Gamma$ ला
  $\TRule{\True}{\lif}$ लावून शाखेचा विस्तार केला आहे: सराव.
\item \Cut{} लावून शाखेचा विस्तार केला आहे. यामुळे दोन शाखा मिळतात: एका
  शाखेत $\sFmla{\True}{\psi}$ आणि दुसऱ्या शाखेत $\sFmla{\False}{\psi}$ असते.
  $\Sat{M}{\Gamma}$ आणि एकतर
  $\Sat{M}{\psi}$ किंवा
  $\Sat/{M}{\psi}$ असल्यामुळे,
  $\Struct{M}$ हे डावीकडील किंवा उजवीकडील शाखा
  पूर्ण करते.
\end{enumerate}
\end{proof}


\begin{prob}
\ref{fol:tab:sou:thm:tableau-soundness} ची सिद्धता पूर्ण करा.
\end{prob}




\begin{cor}
\label{fol:tab:sou:cor:weak-soundness}
जर $\Proves \varphi$, तर $\varphi$ हे वैध आहे.
\end{cor}

\begin{cor}
\label{fol:tab:sou:cor:entailment-soundness}
जर $\Gamma \Proves \varphi$, तर $\Gamma \Entails \varphi$.
\end{cor}

\begin{proof}
  जर $\Gamma \Proves \varphi$, तर काही $\psi_1$, \dots, $\psi_n \in
  \Gamma$ साठी $\{\sFmla{\False}{\varphi}, \sFmla{\True}{\psi_1}, \dots,
  \sFmla{\True}{\psi_n}\}$ चा बंद टॅब्लो असतो.
  \ref{fol:tab:sou:thm:tableau-soundness} नुसार, प्रत्येक
  संरचना~$\Struct{M}$
  काही $\psi_i$ असत्य करते किंवा $\varphi$ सत्य करते. म्हणून
  $\Sat{M}{\Gamma}$ असेल, तर
  $\Sat{M}{\varphi}$ सुद्धा असेल.
\end{proof}

\begin{cor}
\label{fol:tab:sou:cor:consistency-soundness}
जर $\Gamma$ पूर्ततायोग्य असेल, तर तो सुसंगत आहे.
\end{cor}

\begin{proof}
प्रतिलोम विधान सिद्ध करू. $\Gamma$ सुसंगत नाही असे समजा. मग
$\psi_1$, \dots, $\psi_n \in \Gamma$ अशी वाक्ये आणि
$\{\sFmla{\True}{\psi_1}, \dots, \sFmla{\True}{\psi_n}\}$ चा बंद
टॅब्लो असतो. \ref{fol:tab:sou:thm:tableau-soundness} नुसार, प्रत्येक
$i=1$, \dots,~$n$ साठी $\Sat{M}{\psi_i}$ अशी
कोणतीही $\Struct{M}$ नसते. पण मग $\Gamma$
पूर्ततायोग्य नसतो.
\end{proof}












\section{टॅब्लो सह एकरूपता}\label{fol:tab:ide:sec}

टॅब्लोमध्ये एकरूपता असल्यास अतिरिक्त निगमन नियम आवश्यक असतात.
$\eq$ चे नियम पुढीलप्रमाणे आहेत ($t$, $t_1$ आणि $t_2$ ही बंद पदे आहेत):

\begin{defish}
\AxiomC{}
\RightLabel{$\eq$}
\UnaryInfC{\sFmla{\True}{\eq[t][t]}}
\DisplayProof
\hfill
\AxiomC{\sFmla{\True}{\eq[t_1][t_2]}}
\noLine
\UnaryInfC{\sFmla{\True}{\varphi(t_1)}}
\RightLabel{$\TRule{\True}{\eq}$}
\UnaryInfC{\sFmla{\True}{\varphi(t_2)}}
\DisplayProof
\hfill
\AxiomC{\sFmla{\True}{\eq[t_1][t_2]}}
\noLine
\UnaryInfC{\sFmla{\False}{\varphi(t_1)}}
\RightLabel{$\TRule{\False}{\eq}$}
\UnaryInfC{\sFmla{\False}{\varphi(t_2)}}
\DisplayProof
\end{defish}
लक्षात घ्या की इतर सर्व नियमांपेक्षा वेगळे, $\TRule{\True}{\eq}$ आणि
$\TRule{\False}{\eq}$ यांना शाखेवर आधीपासून \emph{दोन} चिन्हांकित
सूत्रे असणे आवश्यक आहे, म्हणजे $\sFmla{\True}{\eq[t_1][t_2]}$
आणि $\sFmla{S}{\varphi(t_1)}$ ही दोन्ही.

\begin{ex}
जर $s$ आणि $t$ ही बंद पदे असतील, तर $\eq[s][t], \varphi(s)
\Proves \varphi(t)$:
\begin{oltableau}
  [\sFmla{\False}{\varphi(t)}, just = \TAss
    [\sFmla{\True}{\eq[s][t]}, just = \TAss
      [\sFmla{\True}{\varphi(s)}, just = \TAss
        [\sFmla{\True}{\varphi(t)}, just={\TRule{\True}{\eq}[2, 3]}, close]
      ]
    ]
  ]
\end{oltableau}
याला एकरूप वस्तूंच्या प्रतिस्थापनाचे तत्त्व किंवा लाइब्निझचा नियम म्हणून
ओळखले जाऊ शकते.

टॅब्लो द्वारे $\eq$ सममित आहे, म्हणजे $\eq[s_1][s_2]
\Proves \eq[s_2][s_1]$, हे सिद्ध होते:
\begin{oltableau}
  [\sFmla{\False}{\eq[s_2][s_1]}, just = \TAss
    [\sFmla{\True}{\eq[s_1][s_2]}, just = \TAss
      [\sFmla{\True}{\eq[s_1][s_1]}, just = {$\eq$}
        [\sFmla{\True}{\eq[s_2][s_1]}, just = {\TRule{\True}{\eq}[2, 3]}, close]
      ]
    ]
  ]
\end{oltableau}
येथे ओळ~$2$ ही $\TRule{\True}{\eq}$ साठीचे पहिले आधारभूत
सूत्र, $\sFmla{\True}{\eq[s_1][s_2]}$, आहे. ओळ~$3$ हे
$\sFmla{\True}{\varphi(s_1)}$ या रूपातील दुसरे आधारविधान आहे---$\varphi(x)$ हे
$\eq[x][s_1]$ आहे असे विचार करा; मग $\varphi(s_1)$ म्हणजे
$\eq[s_1][s_1]$ आणि $\varphi(s_2)$ म्हणजे $\eq[s_2][s_1]$.

तसेच $\eq$ संक्रमक आहे, म्हणजे $\eq[s_1][s_2],
\eq[s_2][s_3] \Proves \eq[s_1][s_3]$, हेही ते सिद्ध करतात:
\begin{oltableau}
  [\sFmla{\False}{\eq[s_1][s_3]}, just = \TAss
    [\sFmla{\True}{\eq[s_1][s_2]}, just = \TAss
      [\sFmla{\True}{\eq[s_2][s_3]}, just = \TAss
        [\sFmla{\True}{\eq[s_1][s_3]}, just = {\TRule{\True}{\eq}[3, 2]}, close]
      ]
    ]
  ]
\end{oltableau}
या टॅब्लोमध्ये $\TRule{\True}{\eq}$ चे पहिले आधारभूत सूत्र
ओळ~$3$ वरील $\sFmla{\True}{\eq[s_2][s_3]}$ आहे ($s_2$ ही~$t_1$ ची,
आणि $s_3$ ही~$t_2$ ची भूमिका बजावतात). $\sFmla{\True}{\varphi(s_2)}$
या रूपातील दुसरे आधारविधान ओळ~$2$ आहे. येथे $\varphi(x)$ हे
$\eq[s_1][x]$ आहे असे विचार करा; त्यामुळे $\varphi(s_2)$ हे
$\eq[s_1][s_2]$ (म्हणजे ओळ~$2$) आणि निष्कर्षातील
सूत्र~$\eq[s_1][s_3]$ हे $\varphi(s_3)$ होते.
\end{ex}

\begin{prob}
पुढील गोष्टींसाठी बंद टॅब्लो द्या:
\begin{enumerate}
\item $\sFmla{\False}{\lforall[x][\lforall[y][((x = y \land \varphi(x))
      \lif \varphi(y))]]}$
\item $\sFmla{\False}{\lexists[x][(\varphi(x) \land
    \lforall[y][(\varphi(y) \lif y = x)])]}$,\\
  $\sFmla{\True}{\lexists[x][\varphi(x)] \land
    \lforall[y][\lforall[z][((\varphi(y) \land \varphi(z)) \lif y = z)]]}$
\end{enumerate}
\end{prob}












\section{एकरूपता सह निर्दोषता}\label{fol:tab:sid:sec}

\begin{prop}
एकरूपतेचे नियम असलेले टॅब्लो निर्दोष आहेत: कोणताही बंद टॅब्लो
  पूर्ततायोग्य नसतो.
\end{prop}

\begin{proof}
मागीलप्रमाणे एवढेच दाखवायचे आहे की टॅब्लोची एखादी शाखा पूर्ततायोग्य
असेल, तर तिला~$\eq$ चा कोणताही नियम लावून मिळणारी शाखासुद्धा पूर्ततायोग्य
असते. शाखेवरील चिन्हांकित सूत्रांचा संच~$\Gamma$ असू द्या आणि
$\Struct{M}$ ही $\Gamma$ पूर्ण करणारी संरचना असू द्या.

शाखेचा~$\eq$ वापरून विस्तार केला आहे, म्हणजे चिन्हांकित
सूत्र~$\sFmla{\True}{\eq[t][t]}$ जोडले आहे असे समजा. स्पष्टपणे
$\Sat{M}{\eq[t][t]}$, म्हणून $\Struct{M}$ हे
$\Gamma \cup \{\sFmla{\True}{\eq[t][t]}\}$ सुद्धा पूर्ण करते.

$\TRule{\True}{\eq}$ वापरून शाखेचा विस्तार केला असेल, तर
सूत्र~$\sFmla{S}{\varphi(t_2)}$ जोडतो; परंतु $\Gamma$ मध्ये
$\sFmla{\True}{\eq[t_1][t_2]}$ आणि $\sFmla{\True}{\varphi(t_1)}$ ही दोन्ही
असतात. म्हणून $\Sat{M}{\eq[t_1][t_2]}$ आणि $\Sat{M}{\varphi(t_1)}$.
$s$ हा $s(x) = \Value{t_1}{M}$ असलेला चल-विन्यास असू द्या.
\ref{fol:syn:ass:prop:sentence-sat-true} नुसार
$\Sat{M}{\varphi(t_1)}[s]$. $\varAssign{s}{s}{x}$ असल्यामुळे,
\ref{fol:syn:ext:prop:ext-formulas} नुसार $\Sat{M}{\varphi(x)}[s]$.
$\Sat{M}{\eq[t_1][t_2]}$ असल्यामुळे
$\Value{t_1}{M} = \Value{t_2}{M}$, आणि म्हणून
$s(x) = \Value{t_2}{M}$. \ref{fol:syn:ext:prop:ext-formulas} पुन्हा
लावल्याने $\Sat{M}{\varphi(t_2)}[s]$ सुद्धा मिळते.
\ref{fol:syn:ass:prop:sentence-sat-true} नुसार
$\Sat{M}{\varphi(t_2)}$. $\TRule{\False}{\eq}$ चे प्रकरण याचप्रमाणे आहे.
\end{proof}



\endgroup
